% Standalone: ./scripts/build_supplement.sh
% Within the original paper checkout, build from its root after main.tex:
% latexmk -pdf -synctex=1 -interaction=nonstopmode -halt-on-error -quiet -outdir=supplementary supplementary/supplementary.tex
\documentclass[11pt,a4paper]{article}
\usepackage[T1]{fontenc}
\usepackage{lmodern}
\usepackage[margin=1in]{geometry}
\usepackage{amsmath,amssymb,amsthm}
\usepackage{xr-hyper}
\usepackage[hidelinks]{hyperref}
% Live references in the paper checkout; nonlinked snapshot when standalone.
\newif\ifmainpapercheckout
\IfFileExists{main.tex}{\mainpapercheckouttrue}{\mainpapercheckoutfalse}
\ifmainpapercheckout
  \externaldocument[main-][nocite]{build/latex/main}[../main.pdf]
  \newcommand{\paperref}[1]{\ref{main-#1}}
  \newcommand{\papereqref}[1]{\eqref{main-#1}}
\else
  % Main-paper numbers for standalone supplementary builds.
% Generated by scripts/sync_paper_assets.py from the in-preparation paper.
% main.tex SHA-256: 5203b97ecb21129d3fe09a3b96d1716c767d7c5839a5865d325a05a63235e071
% References use \ref* and create no missing main.pdf links.
\newlabel{main-sec:barriers}{{3}{6}{}{}{}}
\newlabel{main-sec:one-row-proof}{{5}{8}{}{}{}}
\newlabel{main-sec:matrix-proof}{{4}{7}{}{}{}}
\newlabel{main-eq:spectral-bounds}{{4.1}{7}{}{}{}}
\newlabel{main-eq:shift-event}{{5.5}{10}{}{}{}}

  \newcommand{\paperref}[1]{\ref*{main-#1}}
  \newcommand{\papereqref}[1]{\textup{(\ref*{main-#1})}}
\fi

\newcommand{\Ess}{\mathbf{S}}
\newcommand{\tr}{\operatorname{tr}}
\newcommand{\E}{\mathbb E}
\newcommand{\PP}{\mathbb P}
\renewcommand{\thesection}{S\arabic{section}}
\numberwithin{equation}{section}
\newtheorem{theorem}{Theorem}[section]
\newtheorem{lemma}[theorem]{Lemma}
\newtheorem{proposition}[theorem]{Proposition}
\newtheorem{corollary}[theorem]{Corollary}

\title{Supplementary materials for\\
Sharp tail bounds for the extreme singular values\\
of sub-Gaussian matrices}
\author{Marcus Webb\\University of Manchester}
\date{21 October 2026}
\hypersetup{
  pdftitle={Supplementary materials for Sharp tail bounds for the extreme singular values of sub-Gaussian matrices},
  pdfauthor={Marcus Webb}
}

\begin{document}
\maketitle
\begin{abstract}
This entire document, including its mathematical proof and exposition,
was generated by artificial intelligence without direct input from the author.
Marcus Webb nevertheless takes full responsibility for its content.
The accompanying paper proves singular-value tail bounds with unspecified
universal constants. Here we give the detailed estimates and finite arithmetic
certificates that yield explicit values, with sharper values when $K=1$.
The proof uses the paper's barrier construction and has hybrid verification:
Lean checks the analytic argument
conditional on finite rational certificates, checked separately by exact
Python arithmetic.
\end{abstract}

\section{Explicit constants: an AI-generated proof}
\label{app:constants}

% BEGIN CONSTANTS COMMON
The accompanying paper by Marcus Webb and Morgan Lynch-Davies is available
on arXiv. It presents the probabilistic setting and motivation, comparisons with earlier bounds,
applications to randomized numerical linear algebra, and sharpness examples
and discussion of the correction terms. Its proof explains the barrier method
and establishes the tail bounds with unspecified universal constants.
The paper also states both numerical quartets established below.

This supplement gives the detailed derivation of those quartets, including
the $K=1$ refinement, in Section~\ref{app:constants}; Section~\ref{app:lean}
documents the Lean/Python verification and AI assistance.
We refine the proof in Sections~\paperref{sec:barriers}--\paperref{sec:one-row-proof}
by estimating the mean cost and fluctuations of the row shifts separately.
The theorem statement below is self-contained; its proof uses the paper's
deterministic barrier construction. The constants are not claimed to be optimal.

\begin{theorem}[Explicit singular-value tail bounds]
\label{thm:explicit-constants}
There are positive universal constants $\alpha,\beta,\gamma,\delta$
with the following property.
Let $\Ess\in\mathbb R^{m\times n}$, $m\ge n\ge1$, have independent,
centred, isotropic rows $X_i^T$, meaning
$\mathbb EX_i=0$ and $\mathbb EX_iX_i^T=I_n$.
Suppose that, for some $K\ge1$,
\[
 \mathbb E e^{\lambda\langle X_i,v\rangle}
 \le e^{K^2\lambda^2\|v\|_2^2/2}
 \qquad(\lambda\in\mathbb R,\ v\in\mathbb R^n,\ 1\le i\le m).
\]
Put $L=\log(eK)$, with natural logarithms, and let
$\sigma_{\min}(\Ess)$ and $\sigma_{\max}(\Ess)$ denote the smallest
and largest of its $n$ singular values, including zeros.
Then, for every $t\ge0$,
\begin{align}
 \mathbb P\left\{\sigma_{\min}(\Ess)\le\sqrt m-K\sqrt n
             -\alpha KL-\beta K\sqrt L\,t\right\}
 &\le e^{-t^2/2},\label{eq:supp-lower-tail}\\
 \mathbb P\left\{\sigma_{\max}(\Ess)\ge\sqrt m+K\sqrt n
             +\gamma KL+\delta K\sqrt L\,t\right\}
 &\le e^{-t^2/2}.\label{eq:supp-upper-tail}
\end{align}
For every $K\ge1$, one may take
\begin{equation}\label{eq:supp-general-constants}
 (\alpha,\beta,\gamma,\delta)=(3.84,1.58,2.35,1.5).
\end{equation}
When $K=1$, one may instead take
\begin{equation}\label{eq:supp-K-one-constants}
 (\alpha,\beta,\gamma,\delta)=(1.32,1.12,2.32,\sqrt2).
\end{equation}
\end{theorem}

\paragraph{Guide to the proof.}
The coefficients $\alpha,\gamma$ measure the additive correction $KL$;
$\beta,\delta$ measure the fluctuation scale $K\sqrt L\,t$.
The proof separates these two costs:
\begin{enumerate}
\item Section~\ref{app:transfer} converts one-row Laplace or MGF estimates
into singular-value tails, keeping the correction and fluctuation costs separate.
\item Sections~\ref{app:weighted} and~\ref{app:arithmetic} collect the common
probabilistic tools and the exact arithmetic rules used in the certificates.
\item Sections~\ref{app:low}--\ref{app:low-two-scale} prove the general lower
pair $(\alpha,\beta)$; Section~\ref{app:K-one-proof} sharpens it at $K=1$.
\item Section~\ref{app:upper-constants} proves both upper pairs
$(\gamma,\delta)$, first controlling the mean and then the exponential moment
of the upper shift.
\end{enumerate}
The scalar certificates are included where they enter the proof, so that each
numerical bound can be traced to a finite check. The implementation and its
correspondence with the Lean hypotheses are described in
Sections~\ref{app:lean-certificates}--\ref{app:lean-audit}.

\paragraph{Notation and barrier shifts.}
For symmetric matrices, $R\succ0$ and $R\succeq0$ mean positive definite and
positive semidefinite, respectively. We write $\|R\|$ for the operator norm
and $\|R\|_F=(\operatorname{tr}(R^2))^{1/2}$ for the Frobenius norm.
The matrix $R$ is the current barrier resolvent. The nonnegative exact shifts
$\Delta_-$ and $\Delta_+$ preserve the trace of the resolvent after a row update;
their defining equations are recalled in Sections~\ref{app:low}
and~\ref{app:upper-constants}. In a one-row estimate $R$ is deterministic;
in the matrix construction it is determined by the preceding rows.
The theorem's parameter $t$ is reserved throughout for the final tail bound.

\subsection{From one-row estimates to the four constants}
\label{app:transfer}

\begin{lemma}[Separate drift and fluctuation costs]\label{lem:app-transfer}
Fix $K\ge1$. Suppose $w,\nu>0$ have one of the following properties
for every deterministic $R\succ0$ and every row $X$ satisfying the
assumptions of Theorem~\ref{thm:explicit-constants}.
\begin{enumerate}
\item Whenever $K^2\tr R+w\|R\|_F\le b$, $b>0$, the exact lower shift satisfies
\[
 \E e^{-\lambda(1+b)\Delta_-}
 \le \exp\!\left(-\frac{\lambda}{1+\nu\lambda}\right)
 \qquad(\lambda>0).
\]
Then the lower tail of the theorem holds with correction $w/K$ and
fluctuation term $\sqrt{\nu/2}\,t$.
\item Whenever $K^2\tr R+w\|R\|_F\le b<1$, the exact upper shift satisfies
\[
 \E e^{\lambda(1-b)\Delta_+}
 \le \exp\!\left(\frac{\lambda}{1-\nu\lambda}\right)
 \qquad(0<\lambda<1/\nu).
\]
Then the upper tail of the theorem holds with correction $w/K$ and
fluctuation term $\sqrt{\nu/2}\,t$.
\end{enumerate}
Here $w$ controls the correction cost and $\nu$ controls the fluctuations.
In particular, $w=A K^2L$ and $\nu=B K^2L$ give the
correction coefficient $A$ and fluctuation coefficient $\sqrt{B/2}$.
\end{lemma}

\begin{proof}
Use the same trace-preserving row shifts and outward corrections as
in Section~\paperref{sec:matrix-proof}, now imposing
$K^2\tr R+w\|R\|_F\le b$. Write $S_\pm$ for the sums of the $m$
adaptive row shifts $\Delta_{\pm,i}$ in the two constructions.
Starting at distance
$(K^2n+w\sqrt n)/b$, the trace integral gives
\[
 \sum_i\omega_i\le\frac{w\sqrt n}{b},\qquad
 d_w=\sqrt{K^2n+2w\sqrt n}
 <D_w:=K\sqrt n+\frac wK.
\]
Thus the squared singular values are bounded by the sums of row shifts,
with total initial and outward cost at most $d_w^2/b$.
This is the same calculation as \papereqref{eq:spectral-bounds}, with $w$
in place of $v$ in the regularity condition only.

Take $t>0$ and put $s_\pm=\sqrt m\pm\sqrt{\nu/2}\,t$.
For the lower tail, the target $s_--D_w<0$ gives an empty event.
Otherwise $s_->d_w>0$, so choose
$b=d_w/(s_--d_w)$ and
$\lambda=(\sqrt m-s_-)/(\nu s_-)>0$.
Condition on preceding rows and iterate the lower one-row estimate.
This standard exponential iteration is also used in
\cite[Theorem~2.6]{fan2015exponential}.
Markov's inequality gives
\[
 \PP\{(1+b)S_-\le s_-^2\}
 \le\exp\!\left(\lambda s_-^2-\frac{m\lambda}{1+\nu\lambda}\right)
 =\exp\!\left(-\frac{(\sqrt m-s_-)^2}{\nu}\right)
 =e^{-t^2/2}.
\]
Outside this event,
\[
 \sigma_{\min}(\Ess)^2>\frac{s_-^2}{1+b}-\frac{d_w^2}{b}
 =(s_--d_w)^2>(s_--D_w)^2.
\]
This includes a zero target because $D_w>d_w$.

For the upper tail take $b=d_w/(s_++d_w)<1$ and
$\lambda=(s_+-\sqrt m)/(\nu s_+)$, so $0<\lambda<1/\nu$.
The corresponding iteration and Markov inequality give
\[
 \PP\{(1-b)S_+\ge s_+^2\}
 \le\exp\!\left(-\lambda s_+^2+\frac{m\lambda}{1-\nu\lambda}\right)
 =\exp\!\left(-\frac{(s_+-\sqrt m)^2}{\nu}\right)
 =e^{-t^2/2}.
\]
Outside this event,
\[
 \sigma_{\max}(\Ess)^2<\frac{s_+^2}{1-b}+\frac{d_w^2}{b}
 =(s_++d_w)^2<(s_++D_w)^2.
\]
Taking square roots proves both claims; $t=0$ is immediate.
\end{proof}

\subsection{Gaussian comparison and weighted quadratic forms}
\label{app:weighted}

The following form of the Hsu--Kakade--Zhang comparison keeps the
constants visible. It is \cite[Remark~2.3]{hsu2012tail}; we include
the short deduction for the normalization used here.

\begin{lemma}[Gaussian comparison]\label{app:gaussian-comparison}
For $B\succeq0$ and $\lambda\ge0$ such that $I-2K^2\lambda B\succ0$,
\[
 \E e^{\lambda X^TBX}\le\det(I-2K^2\lambda B)^{-1/2}.
\]
\end{lemma}

\begin{proof}
For a standard Gaussian vector $g$ independent of $X$, the row MGF
assumption gives
\[
 \E e^{\lambda X^TBX}
 =\E_g\E_X e^{\sqrt{2\lambda}\langle B^{1/2}g,X\rangle}
 \le\E_g e^{K^2\lambda g^TBg}
 =\det(I-2K^2\lambda B)^{-1/2}.
\]
\end{proof}

\begin{lemma}[Weighted entropy inequality]\label{app:weighted-entropy}
Let $W\ge0$ have $\E W=1$ and finite entropy
$D(W)=\E(W\log W)$, where $0\log0=0$.
Every random variable $Z$ with $\E e^Z<\infty$ satisfies
\begin{equation}\label{eq:app-entropy-variational}
 \E(WZ)\le D(W)+\log\E e^Z,
\end{equation}
whenever its left side is defined.
No independence between $W$ and $Z$ is required.
\end{lemma}

\begin{proof}
Use the probability measure $W\,d\PP$. Apply concavity of the
logarithm to $e^Z/W$ on $\{W>0\}$, first using bounded truncations
if necessary;
the expectation under $W\,d\PP$ is at most $\E e^Z$.
This proves \eqref{eq:app-entropy-variational}.
\end{proof}


\subsection{Exact arithmetic certificates}
\label{app:arithmetic}

The numerical estimates below use rational enclosures rather than
floating-point approximations. We use two rules: Taylor series with explicit
remainders bound exponentials and logarithms, and Bernstein coefficients
certify polynomial inequalities on whole intervals or rectangles.

\subsubsection{Exponential and logarithmic enclosures}
For integers $N\ge0$, put
$T_N(x)=\sum_{j=0}^{N}x^j/j!$.
Then $T_N(x)\le e^x$ for every $x\ge0$, and
\begin{equation}\label{eq:app-arithmetic-exp}
 T_N(x)\le e^x\le
 T_N(x)+\frac{x^{N+1}}{(N+1)!}\frac1{1-x/(N+2)}
 \qquad(0\le x<N+2).
\end{equation}
For integers $N\ge1$, put
$S_N(z)=2\sum_{j=0}^{N-1}z^{2j+1}/(2j+1)$.
Then
\begin{equation}\label{eq:app-arithmetic-log}
 S_N(z)\le\log\frac{1+z}{1-z}
 \le S_N(z)+\frac{2z^{2N+1}}{(2N+1)(1-z^2)}
 \qquad(0\le z<1).
\end{equation}
Both remainders are bounded by geometric series.
For negative exponentials take reciprocals; for logarithms use
$z=(x-1)/(x+1)$, first reducing $x$ by powers of two when specified.
For $0<x<1$, use $\log x=-\log(1/x)$.
The certificates state their truncation orders and precision:
lower bounds are rounded down and upper bounds up, with each substitution
chosen according to the sign of its multiplier.

\subsubsection{Bernstein coefficients}
Let $p(x)=\sum_{r=0}^d c_rx^r$ and $h>0$. Put
\[
 e_j=\sum_{r=j}^d c_r\binom rj A^{r-j}h^j,\qquad
 b_i=\sum_{j=0}^i e_j\frac{\binom ij}{\binom dj}.
\]
Then
\[
 p(A+hu)=\sum_{i=0}^d b_i\binom di u^i(1-u)^{d-i}.
\]
The identity follows by expanding both sides in powers of $u$.
Its basis functions are nonnegative and sum to one for $0\le u\le1$,
so $p(x)\ge\min_i b_i$ on $[A,A+h]$.
Applying this rule in each coordinate gives the same property for
tensor-product Bernstein coefficients on a rectangle.


% END CONSTANTS COMMON

% BEGIN CONSTANTS LOWER
\subsection{The lower-tail constants}\label{app:low}

Lower-barrier criteria appear in Batson, Spielman and
Srivastava~\cite[Lemma~3.4]{batson2012twice}, Kolla et
al.~\cite[Lemma~3.11]{kolla2010subgraph}, and
Yaskov~\cite[Lemmas~4.1--4.2]{yaskov2014lower} and
\cite[Lemma~4.1]{yaskov2016controlling}. We use the exact trace-preserving
shift and keep its mixed quadratic moment explicit. The proof has three
parts: a comparison for one row, an estimate of its cost in terms of the
spectrum, and a bound on the total cost of restoring the regularity cap.

\subsubsection{Fourth-moment bounds}
Suppose $Z$ is centred, $\mathbb EZ^2=1$, and
$\mathbb E e^{uZ}\le e^{K^2u^2/2}$ for every real $u$.
Write $L=\log(eK)$ and $m_4=\mathbb EZ^4$. Moment interpolation gives
$\mathbb EZ^{2j}\ge m_4^{j-1}$ for every integer $j\ge1$, and therefore
\[
 e^{K^2u^2/2}\ge\mathbb E\cosh(uZ)
 \ge1+\frac{\cosh(u\sqrt{m_4})-1}{m_4}
 \ge\frac{e^{u\sqrt{m_4}}}{2m_4}.
\]
Taking $u=\sqrt{m_4}/K^2$ gives $m_4/(2K^2)\le\log(2m_4)$.
The difference $x/(2K^2)-\log(2x)$ increases for $x>2K^2$.
At $x=5K^2L$ it is $L/2+2-\log(10L)\ge3-\log20>0$, so
\begin{equation}\label{eq:app-low-fourth}
 \mathbb EZ^4\le5K^2L.
\end{equation}
We also need the slightly different bound
\begin{equation}\label{eq:app-low-fourth-offset}
 \mathbb EZ^4\le\frac{499}{100}K^2L+\frac12.
\end{equation}
Put $c=499/100$, $a=c/2-2=99/200$ and $E=e^{2-2L}$.
At $x=cK^2L+1/2$, the same increasing difference is at least
\[
 aL+2-\log(2cL)+E\left(\frac14-\frac1{2cL}\right)
 >aL+2-\log(2cL)+\frac7{50}E.
\]
For $1\le L\le5/2$, this is greater than
$3-\log(1996/99)+(7/50)e^{-3}>-1/250+1/150>0$.
For $L\ge5/2$, the first three terms increase and are already positive
at $5/2$, since $\log(499/20)<259/80$.
The numerical comparisons follow from degree-$24$ Taylor bounds:
$e^{751/250}>1996/99$, $e^3<21$, and $e^{259/80}>499/20$.
This proves \eqref{eq:app-low-fourth-offset}.

For $P\succeq0$, $\operatorname{tr}P=1$, Gaussian averaging gives
\begin{equation}\label{eq:app-low-Tmoment}
 T=X^TPX:\quad \mathbb ET=1,\qquad
 \mathbb ET^2\le\frac53K^2L[1+2\operatorname{tr}(P^2)],\qquad
 \mathbb ET^2\le\frac{499}{100}K^2L+\frac12.
\end{equation}
Indeed $3\mathbb ET^2=\mathbb E_g\mathbb E_X(g^TP^{1/2}X)^4$,
and $\mathbb E_g\|P^{1/2}g\|^4=1+2\operatorname{tr}(P^2)\le3$.

\subsubsection{A quadratic comparison for the exact lower shift}
For $R\succ0$, let $\Delta_-\ge0$ be the unique number such that
$R^{-1}+XX^T-\Delta_-I\succ0$ and
$\operatorname{tr}(R^{-1}+XX^T-\Delta_-I)^{-1}=\operatorname{tr}R$.
Put
\[
 \eta=\|R\|_F,\quad P=R^2/\eta^2,\quad T=X^TPX,
 \quad\kappa=\|R\|/\eta.
\]
We seek a positive matrix $B$ for which
\begin{equation}\label{eq:app-low-comparison}
 T\le(1+X^TBX)\Delta_-.
\end{equation}
To prove this, we may suppose $X\ne0$. Let
$G=(R^{-1}+XX^T)^{-1}$ and
$F=(G^{-1}-\Delta_-I)^{-1}$, and put
\[
 q=\operatorname{tr}(R-G)=\operatorname{tr}(F-G),\qquad
 c_X=\frac{X^TR^3X}{X^TR^2X}.
\]
The rank-one inverse formula and Cauchy--Schwarz give
\[
 q=\frac{X^TR^2X}{1+X^TRX}\le c_X\le\eta,
 \qquad \operatorname{tr}(G^2)=\eta^2-2q c_X+q^2.
\]
Since $F-G=\Delta_-GF\succeq0$,
\[
 \frac q{\Delta_-}=\operatorname{tr}(GF)
 \le\operatorname{tr}(G^2)+q\|G\|.
\]
Using $q\le c_X$ and $\|G\|\le\|R\|$ gives
$q/\Delta_-\le\eta^2+q(\|R\|-c_X)$. Alternatively,
$\|G\|\le\|G\|_F$ gives
\[
 \begin{split}
 \frac q{\Delta_-}
 &\le\eta^2+q\left[q-2c_X+\sqrt{\eta^2-2q c_X+q^2}\right]\\
 &\le\eta^2+q\left[\sqrt{\eta^2-c_X^2}-c_X\right].
 \end{split}
\]
The bracket in the first line increases with $q\in[0,c_X]$.
Thus, with $z=c_X/\eta$,
\begin{equation}\label{eq:app-low-two-comparisons}
 \frac q{\Delta_-}\le\eta^2+q\eta
 \left[\min\{\kappa,\sqrt{1-z^2}\}-z\right].
\end{equation}
In particular, $B_0=R+P(\kappa\eta I-R)$ satisfies
\eqref{eq:app-low-comparison}, by substituting
$q=X^TR^2X/(1+X^TRX)$. More generally,
$B=R+P(A\eta I-cR)$ does so whenever
\[
 \min\{\kappa,\sqrt{1-z^2}\}+(c-1)z\le A\qquad(0\le z\le1).
\]
We will use this freedom when the spectrum is concentrated.

\subsubsection{From the mixed moment to a Laplace bound}
The comparison reduces the probabilistic estimate to two quantities:
the second moment of $T$ and the mixed cost $\mathbb E(TX^TBX)$.
The first determines the fluctuation constant; the second will determine
the regularity condition and hence the correction constant.

Whenever such a matrix also satisfies $\mathbb E(TX^TBX)\le b$,
let $m_2\ge\mathbb ET^2$. Weighted Cauchy--Schwarz, as in
\cite[Lemma~4.2]{yaskov2014lower}, and $e^{-x}\le(1+x)^{-1}$ give
\begin{equation}\label{eq:app-lower-moment-laplace}
 \begin{aligned}
 \mathbb Ee^{-\lambda(1+b)\Delta_-}
 &\le1-\lambda(1+b)\mathbb E
   \frac{T}{1+X^TBX+\lambda(1+b)T}\\
 &\le1-\frac{\lambda}{1+m_2\lambda}
 \le\exp\left(-\frac{\lambda}{1+(m_2-1/2)\lambda}\right)
 \qquad(\lambda>0).
 \end{aligned}
\end{equation}
The last inequality follows by putting
$z=\lambda/[2+(2m_2-1)\lambda]$ in
$\log[(1+z)/(1-z)]\ge2z$.
Taking $m_2=(499/100)K^2L+1/2$, every mixed-cost cap proved below gives
$\nu_-=(499/100)K^2L$, and hence $\beta=79/50$, since
$499/100<2(79/50)^2$.

\subsection{Spectral costs and the lower correction budget}
\label{app:low-two-scale}

It remains to control the mixed cost without paying the worst spectral cost
at every step of the barrier construction. We first obtain bounds depending
on the spectral shape $\kappa=\|R\|/\|R\|_F$, then integrate that cost
along the outward corrections. This is the step that yields $\alpha=3.84$.

\subsubsection{A determinant envelope}
Normalize $\eta=1$. The eigenvalue of $B_0$ corresponding to an
eigenvalue $\rho$ of $R$ is $b_0(\rho)=\rho+\rho^2(\kappa-\rho)$.
For $0<q<1$ put
\[
 \mathcal J(q)=-\log(1-q)-q,\qquad
 \mathcal M(q)=q-\frac{2(1-q)}q\mathcal J(q)
     =\sum_{j=2}^\infty\frac{2q^j}{j(j+1)}.
\]
If $q=u\kappa$ and $\mathcal M(q)\ge\kappa^2$, then
\begin{equation}\label{eq:app-low-shape-envelope}
 -\log[1-u b_0(\rho)]-u\rho
 \le\frac{\rho^2}{\kappa^2}\mathcal J(u\kappa),\qquad0\le \rho\le\kappa.
\end{equation}
Here $\kappa-b_0(\rho)=(\kappa-\rho)(1-\rho^2)\ge0$, so the logarithm is well-defined.
To prove the inequality in the strict case, set
$c_0=\mathcal J(q)/\kappa^2$, $d(\rho)=1-u b_0(\rho)$, and
$\mathcal F(\rho)=u\rho+c_0\rho^2+\log d(\rho)$.
Then $\mathcal F(0)=\mathcal F(\kappa)=0$ and
\[
 \begin{split}
 d(\rho)\mathcal F'(\rho)/\rho={}&(2c_0-u^2-2u\kappa)
 +(3u-u^2\kappa-2uc_0)\rho\\
 &+(u^2-2uc_0\kappa)\rho^2+2uc_0\rho^3.
 \end{split}
\]
This cubic is positive at zero and negative at $\kappa$.
Indeed the latter sign is precisely $\mathcal M(q)>\kappa^2$, and the former
follows from
$\mathcal M(q)<\mathcal J(q)/q-q/2$ by comparing their positive series.
Its positive leading coefficient forces one root in $(0,\kappa)$ and
another in $(\kappa,\infty)$; the third root is negative, since their
product is negative. Thus $\mathcal F$ first increases and then decreases on
$[0,\kappa]$, proving $\mathcal F\ge0$. The non-strict case follows by continuity.

\subsubsection{A uniform bound for concentrated spectra}
When $\kappa$ is large we use a comparison matrix adapted to that range,
giving a constant upper bound for the mixed cost.

\begin{lemma}[A uniform cost for concentrated spectra]
\label{lem:general-lower-plateau}
Let $\|R\|_F=1$, $\kappa=\|R\|\ge71/100$, $P=R^2$ and
$T=X^TPX$. There is a positive matrix $B$ satisfying the lower-shift
comparison $T\le(1+X^TBX)\Delta_-$ and
\[
 \mathbb E(TX^TBX)-K^2\operatorname{tr}R
       \le\frac{4137}{1000}K^2L.
\]
\end{lemma}
\begin{proof}
First suppose $71/100\le\kappa\le147/200$, and set
\[
 s=\frac{347}{500},\quad u=\frac{1207}{1000},\quad
 w=\frac{26}{25},\quad c=\frac{17}{10},\qquad
 A(k)=k+\frac7{10}\frac{1+s^2-k^2}{2s}.
\]
Take $B=R+P(A(\kappa)I-cR)$.
The tangent inequality $\sqrt{1-k^2}\le(1+s^2-k^2)/(2s)$ and
\eqref{eq:app-low-two-comparisons} give the required shift comparison,
since
\[
 \min\{\kappa,\sqrt{1-z^2}\}+\frac7{10}z
 \le\kappa+\frac7{10}\sqrt{1-\kappa^2}\le A(\kappa).
\]
For the first inequality split at $z=\sqrt{1-\kappa^2}$.
On the left use $\kappa+(7/10)z$; on the right the function
$\sqrt{1-z^2}+(7/10)z$ decreases, since $\kappa^2\le100/149$.

We next give a scalar determinant majorant valid throughout this range.
Set
\[
 \begin{gathered}
 A_0=\frac65,\quad A_1=\frac{1211}{1000},\quad
 v_0=\frac{1631}{1000},\quad v_1=\frac{167}{100},\\
 v(A)=v_0+\frac{v_1-v_0}{A_1-A_0}(A-A_0).
 \end{gathered}
\]
The increasing polynomial $A(k)$ takes values in $[A_0,A_1]$ for
$71/100\le k\le147/200$. For $A\in[A_0,A_1]$, define
\[
 \begin{aligned}
 d_A(\rho)&=1-u(\rho+A\rho^2-c\rho^3),\\
 \psi_A(\rho)&=u\rho+2v(A)\rho^2-2w\rho^4,\\
 g_A(\rho)&=\psi_A(\rho)+\log d_A(\rho).
 \end{aligned}
\]
We claim that
\begin{equation}\label{eq:general-plateau-scalar}
 d_A(\rho)>0,\qquad g_A(\rho)\ge0,
       \qquad 0\le\rho\le147/200.
\end{equation}
It suffices to check $A=A_0,A_1$: $d_A$ is affine in $A$, and
$g_A$ is concave in $A$. Positivity of $B$ follows at the same time
from $1+A\rho-c\rho^2\ge1-c(147/200)^2>0$.

Here are exact certificates for the two endpoints. The derivative of
$e^{\psi_A}d_A$ is $\rho e^{\psi_A}p_A(\rho)$, where the coefficients
of $p_A$, in increasing order, are
\[
 \begin{gathered}
 4v-u^2-2uA,\quad 3uc-u^2A-4uv,\quad u^2c-4uvA-8w,\\
 4uvc+8uw,\quad 8uwA,\quad -8uwc.
 \end{gathered}
\]
Write $b=147/200$. For both endpoints, the coefficients of
$(1+x)^5p_A(bx/(1+x))$ have signs $+,+,+,-,-,+$.
Moreover $p_A(0)>0$, and its second root lies in $[r_-,r_+]$ with
integer endpoint numerators, over the common denominator $10^5$,
\[
 (63549,63550)\quad\text{for }A_0,
 \qquad(64243,64244)\quad\text{for }A_1.
\]
Direct substitution gives $p_A(r_-)<0<p_A(r_+)$ and
\[
 d_A(r_+)T_{24}(\psi_A(r_-))>1.
\]
The additional rational checks
$1+2Ab-3cb^2>0$, $d_A(b)>0$ and
$u+4vr_--8wr_+^3>0$ show that $d_A$ decreases and is positive on
$[0,b]$, while $\psi_A$ increases on the root bracket.
Descartes' rule therefore gives exactly two critical points, and the
minimum of $e^{\psi_A}d_A$ is attained at zero or at the second one.
This proves \eqref{eq:general-plateau-scalar}.

Writing $\zeta=\operatorname{tr}(P^2)$, Gaussian comparison, with one
largest eigenvalue kept exactly, now gives
\[
 \log\mathbb E\exp\left\{\frac{u}{2K^2}
       (X^TBX-K^2\operatorname{tr}R)\right\}
 \le v(A(\kappa))-w\zeta-\frac12g_{A(\kappa)}(\kappa).
\]
The fourth-moment bound gives
$\mathbb E(T\log T)\le\log[5K^2L(1+2\zeta)/3]$.
Maximizing $-w\zeta+\log[(1+2\zeta)/3]$ and then using
$\sup_{L>0}(\log L+C)/L=e^{C-1}$ bounds the mixed cost by $K^2L$ times
\begin{equation}\label{eq:general-plateau-cost}
 \frac4u+\frac{20}{3uw}\frac{e^{E(\kappa)}}{\sqrt{d(\kappa)}},
\end{equation}
where the explicit quartic polynomials are
\[
 \begin{split}
 E(k)&=v(A(k))(1-k^2)-4+\frac w2-\frac{uk}{2}+wk^4,\\
 d(k)&=1-u(k+A(k)k^2-ck^3).
 \end{split}
\]
The logarithmic derivative of the second term in
\eqref{eq:general-plateau-cost} has the sign of
\[Q(k)=2d(k)E'(k)-d'(k).\]
Every Bernstein coefficient of $Q'$ on $[71/100,147/200]$ is less
than $-5$; those of $E'$ and $d'$ are negative.
Moreover $Q(k_-^*)>0>Q(k_+^*)$ for
\[
 k_-^*=\frac{72167}{10^5},\qquad
 k_+^*=\frac{72168}{10^5}.
\]
Thus the maximum is attained between these two points. Since $E$ and
$d$ decrease, its value is bounded by replacing $E(\kappa)$ by $E(k_-^*)$
and $d(\kappa)$ by $d(k_+^*)$. Finally, with $C=4137/1000$,
\[
 T_{24}(-2E(k_-^*))
 >\frac{400}{9w^2(uC-4)^2d(k_+^*)}.
\]
This rational certificate proves that \eqref{eq:general-plateau-cost}
is less than $C$.

For $\kappa\ge147/200$, use $B_2=R+P((74/55)I-(19/10)R)$.
It is positive, and it satisfies \eqref{eq:app-low-comparison}, because
$\sqrt{1-z^2}+(9/10)z\le\sqrt{181}/10<74/55$.
Take
\[
 a=\frac{147}{200},\quad A=\frac{74}{55},\quad c=\frac{19}{10},\quad
 u=\frac65,\quad v=\frac{27}{16},\quad w=\frac{21}{20},
\]
Define $d_2(\rho)=1-u(\rho+A\rho^2-c\rho^3)$,
$\psi_2(\rho)=u\rho+2v\rho^2-2w\rho^4$, and $g_2=\psi_2+\log d_2$.
The same derivative formula and Descartes test on $[0,1]$ have signs
$+,+,-,-,+,+$, second-root bracket $[607/1000,608/1000]$, and
\[
 d_2(608/1000)T_{24}(\psi_2(607/1000))>1.
\]
On that bracket $d_2$ decreases and $\psi_2$ increases.
Thus $d_2>0$, $g_2\ge0$ on $[0,1]$, and $g_2$ increases above $608/1000$.
Keeping its contribution at the largest eigenvalue therefore gives
\eqref{eq:general-plateau-cost} with
$d=d_2(a)$ and $E=v-4+w/2-ua/2-va^2+wa^4$.
The certificate
\[
 T_{24}(-2E)
 >\frac{400}{9w^2(uC-4)^2d}
\]
proves the bound on this final range as well.
\end{proof}

\subsubsection{A spectral cost profile from entropy}
For the remaining shapes, an entropy estimate gives a smaller cost that
depends on $\kappa$. We combine it with the preceding uniform bound in
the following profile.

Normalize $\|R\|_F=1$, and retain $P=R^2$, $T=X^TPX$,
$\kappa=\|R\|$ and $B_0=R+P(\kappa I-R)$. Define
\begin{equation}\label{eq:app-new-lower-profile}
 h_*(k)=\min\left\{\frac{4137}{1000},
    \max\left\{\frac32,\frac{57}{40}+\frac{23}{10}k
                    +\frac{47}{20}k^2-\frac3{10}k^3\right\}\right\}.
\end{equation}
We will prove the mixed-cost estimate
\begin{equation}\label{eq:app-new-lower-mixed}
 \mathbb E(TX^TB_0X)-K^2\operatorname{tr}R
       \le K^2L h_*(\kappa)
\end{equation}
whenever $h_*(\kappa)<4137/1000$.
Together with Lemma~\ref{lem:general-lower-plateau}, this gives the lower
mixed-cost bound for every spectral shape.

\paragraph{The entropy estimate.}
For every real $p\ge19/4$, put
$\varepsilon_p=-\log(1-e^{-(p-1/20)})$. Then
\begin{equation}\label{eq:app-moment-exponential}
 \sup_{y\ge0}\frac{y^p}{e^y-1}
 \le \frac{p^p e^{-p}}{1-e^{-(p-1/20)}}.
\end{equation}
Indeed the unique positive maximizer $y_p$ satisfies
$y_p=p(1-e^{-y_p})$. The inequality
$e^{p-1/20}>20p$, valid for $p\ge19/4$, gives $y_p>p-1/20$;
now use $y^pe^{-y}\le p^pe^{-p}$.
The initial exponential inequality follows from
$T_{24}(47/10)>95$ and the monotonicity of $e^p/p$.

The free moment order $p$ lets us balance the entropy cost against the
Gaussian determinant bound. Choose $0<q<1$, put $\mathsf B=2\kappa^2/q$, and write
$M=(1-q)^{-1/(2\kappa^2)}$.
Gaussian comparison and $\operatorname{tr}P=1$, $\|P\|=\kappa^2$
give $\mathbb E e^{T/(K^2\mathsf B)}\le M$.
Applying \eqref{eq:app-moment-exponential} pointwise and then Jensen's
inequality under the probability measure with density $T$ gives
\begin{equation}\label{eq:app-new-entropy-bound}
 \begin{split}
 \mathbb E(T\log T)
 &\le\frac{\log\mathbb ET^p}{p-1}\\
 &\le\frac{2p(L-1)+p\log(\mathsf Bp)-p
                   +\log(M-1)+\varepsilon_p}{p-1}.
 \end{split}
\end{equation}
Consequently, if $u\kappa<1$, $\mathcal M(u\kappa)\ge\kappa^2$, and
\begin{equation}\label{eq:app-new-entropy-condition}
 (p-1)\frac{\mathcal J(u\kappa)}{2\kappa^2}
       +\log(M-1)+p\log(\mathsf Bp)+\varepsilon_p\le3p,
\end{equation}
then \eqref{eq:app-low-shape-envelope}, the entropy variational
inequality, and \eqref{eq:app-new-entropy-bound} imply
\[
 \mathbb E(TX^TB_0X)-K^2\operatorname{tr}R
 \le \frac{4p}{(p-1)u}K^2L.
\]
Thus the choice $u=4p/[(p-1)h_*(\kappa)]$ proves
\eqref{eq:app-new-lower-mixed} once the displayed scalar conditions hold.

\subsubsection{Certifying the profile on the full spectral range}
We now verify the scalar conditions for the chosen profile. The interval
checks cover the middle range of $\kappa$; a direct bound covers small
$\kappa$, and the uniform bound above covers the remaining concentrated
spectra.

For $3/100\le\kappa\le711/1000$, use the cells
$I_j=[j/1000,(j+1)/1000]$, $30\le j<711$.
Writing $\bar k_j=(2j+1)/2000$ for the midpoint, choose the constant power
\[
 p_j=\begin{cases}
 \dfrac{37}{10}+\dfrac{13}{10\bar k_j}-\dfrac{\bar k_j^2}{2},&30\le j<650,\\[4pt]
 12-10\bar k_j,&650\le j<711.
 \end{cases}
\]
Every power is at least $19/4$.
Put
\[
 D_p(k)=1+2pk^2-4k^4,\quad
 q_p(k)=1-D_p(k)^{-1},\quad
 \mathsf B_p(k)=2k^2+\frac1{p-2k^2}.
\]
These choices satisfy $\mathsf B_p=2k^2/q_p$ and
$\log M_p(k)=\log D_p(k)/(2k^2)$.
The function $\mathsf B_p$ is increasing, and $\log M_p$ is decreasing:
for the latter, $y\mapsto\log(1+2py-4y^2)$ is concave and vanishes at
zero. Also $\mathcal J(uk)/(2k^2)$ and
$\mathcal M(uk)/k^2$ are increasing in both positive variables $u,k$,
as follows from their positive power series.

For $x\ge1$, let $U(x)$ be the logarithm upper bound in
\eqref{eq:app-arithmetic-log} with $N=24$ and $z=(x-1)/(x+1)$,
rounded upwards to a multiple of $10^{-12}$.
On a cell $[a,b]$, let
\[
 u_+=\frac{4p}{(p-1)h_*(a)},\quad
 u_-=\frac{4p}{(p-1)h_*(b)},\quad
 \ell=\frac{U(D_p(a))}{2a^2},
\]
rounding $\ell$ upwards in the same manner. If $\ell\ge8$, set
$V(\ell)=\ell$. Otherwise let $v$ be the reciprocal of the exponential
upper bound in \eqref{eq:app-arithmetic-exp}, with $N=40$ and $x=\ell$,
rounded downwards to a multiple of $10^{-12}$, and put
\[
 V(\ell)=\ell-\sum_{r=1}^{6}\frac{v^r}{r}.
\]
Then $V(\ell)\ge\log(M_p(a)-1)$.
On every one of the $681$ cells, exact rational arithmetic gives
\[
 u_+b<1,\qquad
 \sum_{r=2}^{60}\frac{2(u_-a)^r}{r(r+1)}\ge a^2,
\]
\[
 \begin{aligned}
 &(p-1)\frac{U((1-u_+b)^{-1})-u_+b}{2b^2}
      +V(\ell)+pU(p\mathsf B_p(b))\\
 &\hspace{3em}+\frac1{T_{24}(p-1/20)-1}<3p-\frac1{10000}.
 \end{aligned}
\]
The monotonicities above prove the required inequalities throughout
each cell. These are finite rational checks, with all functions,
subdivisions and powers specified explicitly.

For $0<\kappa\le3/100$, take $p=50$, $u=400/147$, and use the same
entropy and determinant bounds at the larger effective shape
$k=3/100$. Specifically, take $\mathsf B=\mathsf B_{50}(3/100)$;
the bound $\|P\|\le(3/100)^2$ gives the same Gaussian MGF bound,
and $\mathcal J(u\kappa)/(2\kappa^2)$ is bounded by its value at $k$.
Their certificate has the same strict margin $1/10000$.
Here $u^2>3$, so
$\mathcal M(u\kappa)\ge u^2\kappa^2/3>\kappa^2$.
This proves the remaining part of \eqref{eq:app-new-lower-mixed}.
Finally the cubic in \eqref{eq:app-new-lower-profile} exceeds
$4137/1000$ at $711/1000$, so Lemma~\ref{lem:general-lower-plateau} covers all remaining
shapes.

\subsubsection{Integrating a spectral cost along the barrier}
The next lemma converts a shape-dependent regularity condition into a
total initial and outward cost. Its integral replaces a uniform bound
for the profile, which would give a larger correction.

\begin{lemma}[Cost of a mixed regularity condition]\label{lem:app-mixed-cap-cost}
Let $a,b>0$, and let $h:[0,1]\to[0,\infty)$ be nondecreasing and
absolutely continuous, with $h(k)>0$ for $k>0$ and
$2\int_0^1h(k)\,dk\ge h(1)$.
A trace-preserving barrier construction imposing
\[
 a\operatorname{tr}R+\|R\|_Fh(\|R\|/\|R\|_F)\le b
\]
has total initial and outward cost at most
\[
 \frac{an+2\mathcal H\sqrt n}{b},\qquad
 \mathcal H=\int_0^1\frac{k h(k)}{\sqrt{1-k^2}}\,dk,
\]
when started at distance
$[an+\sqrt n\,h(n^{-1/2})]/b$.
In particular, for $h(k)=f+o\min\{k,k_0\}$, where
$f,o\ge0$, $0\le k_0\le1$ and $f+ok_0>0$, one has
\[
 \mathcal H=f+o\mathcal A(k_0),\qquad
 \mathcal A(k_0)=\frac{\arcsin k_0+k_0\sqrt{1-k_0^2}}2.
\]
\end{lemma}

\begin{proof}
For $n>1$, put $\tau=\operatorname{tr}R$, $F=\|R\|_F$, $M=\|R\|$.
An outward shift has $R(\omega)=(R^{-1}+\omega I)^{-1}$.
Both $F$ and $M/F$ decrease, since
\[
 \frac{d}{d\omega}\log\frac MF
 =-M+\frac{\operatorname{tr}(R^3)}{\operatorname{tr}(R^2)}\le0.
\]
Thus the cap decreases continuously to zero and can always be restored.
Cauchy--Schwarz on all eigenvalues except the largest gives
\[
 M\le\frac\tau n+\sqrt{1-\frac1n}
                         \sqrt{F^2-\frac{\tau^2}{n}}.
\]
For $n^{-1/2}\le x\le1$, define
\[
 \psi_n(x)=\frac1{nx}+\sqrt{1-\frac1n}
                              \sqrt{1-\frac1{nx^2}},\qquad
 g(x)=a+xh(\psi_n(x)).
\]
The function $\psi_n$ is nondecreasing, and $g$ is strictly increasing.
Until a correction stops, $b<\tau g(F/\tau)$.
Initially $\tau=b/g(n^{-1/2})$, and the cap already holds whenever
$\tau\le b/g(1)$.
Since $-\tau'=F^2$ and the trace intervals of the corrections are disjoint,
the sum of outward shifts is at most
\[
 \int_{b/g(1)}^{b/g(n^{-1/2})}
 \frac{d\tau}{\tau^2[g^{-1}(b/\tau)]^2}
 =\frac1b\int_{n^{-1/2}}^1\frac{g'(x)}{x^2}\,dx.
\]
Adding the initial distance cancels the lower endpoint after integration
by parts. With $\varphi=\arcsin(n^{-1/2})$ and
$x=(\sqrt n\cos\theta)^{-1}$, the contribution beyond $an/b$,
divided by $\sqrt n/b$, becomes
\[
 I_h(\varphi)=h(1)\sin\varphi+
 2\int_0^{\pi/2-\varphi}
                  h(\sin(\theta+\varphi))\sin\theta\,d\theta.
\]
As $h'\ge0$ almost everywhere,
\[
 \begin{split}
 I_h'(\varphi)
 &=-h(1)\cos\varphi+
 2\int_0^{\pi/2-\varphi}
 h'(\sin(\theta+\varphi))\cos(\theta+\varphi)\sin\theta\,d\theta\\
 &\le\cos\varphi\left[-h(1)+2\int_0^1k h'(k)\,dk\right]
 =\cos\varphi\left[h(1)-2\int_0^1h(k)\,dk\right]\le0.
 \end{split}
\]
The inequality uses $\sin\theta\le\cos\varphi\sin(\theta+\varphi)$
and the substitution $k=\sin(\theta+\varphi)$.
Hence $I_h(\varphi)\le I_h(0)=2\mathcal H$.
For $n=1$ there are no corrections, and the same formula follows directly.
The stated special case satisfies the integral condition, and its
integral is evaluated by splitting at $k_0$.
\end{proof}


\subsubsection{The general lower-tail constants}
We apply the budget lemma to the certified profile and bound its integral
by exact scalar comparisons.

The profile $h_*$ is nondecreasing and absolutely continuous.
Moreover $h_*\ge3/2$ on $[0,1]$ and $h_*\ge3$ on $[1/2,1]$, so
$2\int_0^1h_*(k)\,dk>h_*(1)$.
Restoring the scale $\eta$ gives a valid lower regularity cap
\[
 K^2\operatorname{tr}R+K^2L\|R\|_F
       h_*(\|R\|/\|R\|_F)\le b.
\]
Lemma~\ref{lem:app-mixed-cap-cost} bounds its correction coefficient by
$\mathcal H=\int_0^1kh_*(k)/\sqrt{1-k^2}\,dk$.
For an exact bound, write $\Pi(k)=a_0+a_1k+a_2k^2+a_3k^3$ for the
cubic in \eqref{eq:app-new-lower-profile}, and put $s=713/1000$.
On $[0,s]$, use
$h_*(k)\le\Pi(k)+(3/40)\mathbf1_{\{k\le1/30\}}$;
on $[s,1]$, use $h_*\le4137/1000$.
Integration gives
\[
 \mathcal H\le a_0+\frac{2a_2}{3}+\frac{83}{80}\arcsin s
 +Q\sqrt{1-s^2}+\frac3{40}\left(1-\sqrt{1-(1/30)^2}\right),
\]
where
\[
 Q=\frac{4137}{1000}-a_0-\frac{2a_2}{3}-\frac{a_2s^2}{3}
       -\frac{a_1s}{2}-\frac{a_3s(3+2s^2)}8>0.
\]
The degree-seven alternating sine bound and two squared comparisons give
\[
 \arcsin s<\frac{793769}{10^6},\qquad
 \sqrt{1-s^2}<\frac{701165}{10^6},\qquad
 \sqrt{1-(1/30)^2}>\frac{999444}{10^6}.
\]
Substitution yields
\[
 \mathcal H<\frac{96}{25}.
\]
The transfer argument of Lemma~\ref{lem:app-transfer}, with dimension
cost $K^2n+2K^2L\mathcal H\sqrt n$, now proves the lower tail with
\[
 \alpha=\frac{96}{25}=3.84,\qquad\beta=\frac{79}{50}=1.58.
\]
As before, the strict dimension estimate includes a zero target threshold.
% END CONSTANTS LOWER

% BEGIN K1 LOWER REFINEMENT
\subsection{The lower-tail refinement when \texorpdfstring{$K=1$}{K=1}}
\label{app:K-one-proof}
We improve the mixed quadratic moment in the lower-shift comparison.
Put $\eta=\|R\|_F$, $P=R^2/\eta^2$, $T=X^TPX$ and
$\kappa=\|R\|/\eta$. The exact second moment and the fourth-order
consequence of the directional MGF assumption suffice for the rounded constants.

\subsubsection{A mixed covariance bound}
Compare the fourth-order coefficients in
$\mathbb E\cosh(u\langle X,v\rangle)\le e^{u^2\|v\|^2/2}$,
using its exact second moment, to obtain
$\mathbb E\langle X,v\rangle^4\le3\|v\|^4$.
This is the fourth-cumulant necessary condition for strict
sub-Gaussianity; see \cite[Proposition~3.2]{arbel2020strict}.
The Taylor argument applies here because the directional MGF assumption
ensures all required moments.
Gaussian averaging, as in \eqref{eq:app-low-Tmoment}, therefore gives
\[
 \mathbb ET^2\le1+2\operatorname{tr}(P^2)\le3,
 \qquad
 \operatorname{Var}(X^TAX)\le2\|A\|_F^2\quad(A\succeq0).
\]
For positive semidefinite $A,B$, this also gives
\begin{equation}\label{eq:app-K-one-covariance}
 \operatorname{Cov}(X^TAX,X^TBX)
 \le\|A\|_F\|B\|_F+\operatorname{tr}(AB).
\end{equation}
Indeed, nonnegative variance implies
$4\operatorname{Cov}(U,V)\le\operatorname{Var}(sU+s^{-1}V)$
for every $s>0$. Apply the variance bound to
$sA+s^{-1}B\succeq0$ and minimize over $s$.
Zero matrices follow by continuity.
This reasoning does not assume independence of the coordinates.

\subsubsection{An orthogonal rank-one refinement}
Suppose $P=a uu^T+P_{\perp}$ and $B=d uu^T+B_{\perp}$, where $u$ is a
unit vector, $a,d>0$, and $P_{\perp},B_{\perp}\succeq0$ annihilate $u$.
Put $p=\|P_{\perp}\|_F$, $q=\|B_{\perp}\|_F$. If $p,q>0$ and
\begin{equation}\label{eq:app-K-one-rankone-ratio}
 r=\frac{aq}{dp},\qquad r+r^{-1}\le6,
\end{equation}
then
\begin{equation}\label{eq:app-K-one-rankone-covariance}
 \operatorname{Cov}(X^TPX,X^TBX)\le2ad+2pq.
\end{equation}
Here $a,d$ denote the displayed eigenvalues only.
To prove this, put $Z=\langle X,u\rangle^2$, $v=\operatorname{Var}Z$,
$Y=X^TP_{\perp}X$, $W=X^TB_{\perp}X$, $x=\operatorname{Cov}(Z,Y)$
and $y=\operatorname{Cov}(Z,W)$, so $0\le v\le2$.
Averaging the directional fourth-moment bound at
$\sqrt h\,u+\sqrt j\,g$, with $g\sim N(0,P_{\perp})$, gives
\[
 (2-v)h^2-6xhj+3(2p^2-\operatorname{Var}Y)j^2\ge0
 \qquad(h,j\ge0).
\]
Thus, writing $x_+=\max\{x,0\}$ and similarly for $y_+$,
\[
 \operatorname{Var}Y\le2p^2-\frac{3x_+^2}{2-v},\qquad
 \operatorname{Var}W\le2q^2-\frac{3y_+^2}{2-v}
 \qquad(0<v<2).
\]
If $x,y<0$, ordinary Cauchy--Schwarz proves the claim immediately.
Otherwise $xy\le x_+y_+$. Positivity of the covariance matrix gives
\[
 \operatorname{Cov}(Y,W)\le\frac{xy}{v}
 +\sqrt{(\operatorname{Var}Y-x^2/v)(\operatorname{Var}W-y^2/v)}.
\]
Using the two variance bounds and weighted arithmetic--geometric mean
in the square root, with $\Lambda_v=(1+v)/[v(2-v)]$, bounds the full covariance by
\[
 adv+2pq+dx_++ay_++\frac{x_+y_+}{v}
       -\Lambda_v\left(\frac qp x_+^2+\frac pq y_+^2\right).
\]
Completing the square, the maximum of the last four terms over real
$x_+,y_+$ is
\[
 ad(2-v)\frac{(1+v)(r+r^{-1})+(2-v)}{3(v+4)}
 \le ad(2-v).
\]
This proves the claim for $0<v<2$. At $v=0$, $x=y=0$; at $v=2$ the
fourth-moment inequalities give $x,y\le0$. Cauchy--Schwarz proves
both endpoint cases directly.

\subsubsection{The cost when the spectrum is spread out}
Suppose first that $\kappa\le73/100$. Use the comparison
$T\le(1+X^TB_0X)\Delta_-$ from Section~\ref{app:low-two-scale},
where $B_0=R+P(\|R\|I-R)$.
Normalize $\eta=\|R\|_F=1$, write $\rho_j$ for the eigenvalues of $R$,
and put $\kappa=\max_j\rho_j$ and $s_\ell=\sum_i\rho_i^\ell$, so $s_2=1$.
Bound \eqref{eq:app-K-one-covariance} gives
\[
 \mathbb E(TX^TB_0X)-\operatorname{tr}R
 \le\kappa+\kappa s_4-s_5+\sqrt{s_4}\,\|B_0\|_F
 \le2\kappa.
\]
Here is the spectral inequality needed for the last step.
The number $\kappa-\kappa s_4+s_5$ is positive, and direct expansion gives
\[
 \begin{split}
 &(\kappa-\kappa s_4+s_5)^2-s_4\|B_0\|_F^2\\
 &\quad=(\kappa^2-s_4)(1-2s_4)
 +\sum_{i<j}\rho_i^2\rho_j^2(\rho_i-\rho_j)^2
       [2\kappa(\rho_i+\rho_j)-\rho_i^2\rho_j^2].
 \end{split}
\]
Every bracket is nonnegative. If $s_4\le1/2$, the whole expression is
nonnegative, since $s_4\le\kappa^2$.
Otherwise put $v=\kappa^2>1/2$ and $r_*=\sqrt{1-v}$.
The largest eigenvalue is unique, all other eigenvalues $r$ are at most
$r_*$, and $2s_4-1\le(2v-1)^2$.
The pair involving $\kappa$ and $r$, divided by $r^2(v-r^2)$, equals
\[
 G(r)=v(\kappa-r)
       \left[2\kappa-\frac{vr^2}{\kappa+r}\right].
\]
Here $\kappa-r\ge(2v-1)/(2\kappa)$ and
$vr^2/(\kappa+r)\le\kappa/2$, because $r^2\le1-v<1/2$.
Hence $G(r)\ge(3v/4)(2v-1)\ge(2v-1)^2$, since
$v\le(73/100)^2<4/5$.
Finally, $\sum_{r\ne\kappa}r^2(v-r^2)=v-s_4$.
Thus these pairs alone offset the negative first term.
Taking square roots proves the cost $2\kappa$ in this range.
The next argument covers $\kappa\ge73/100$.

\subsubsection{A uniform cost for concentrated spectra}
We next prove that, when $\kappa\ge73/100$, a positive comparison matrix
$B$ can be chosen with
\begin{equation}\label{eq:app-K-one-uniform-cost}
 T\le(1+X^TBX)\Delta_-,\qquad
 \mathbb E(TX^TBX)\le\operatorname{tr}R+\frac{73}{50}.
\end{equation}
For $c\ge1$ and $A\ge\sqrt{1+(c-1)^2}$, the tangent inequality
$\sqrt{1-z^2}-z\le A-cz$ gives the comparison
$B=R+P(AI-cR)$.
It is positive when $1+A-c>0$, by concavity of $1+Ax-cx^2$ on $[0,1]$.
Convex combinations with $B_0$ remain valid comparison matrices.
Bound \eqref{eq:app-K-one-covariance} and Young's inequality give, for $\xi>0$,
\[
 \mathbb E(TX^TBX)-\operatorname{tr}R
 \le A+\frac1{2\xi}+\sum_i\rho_i^2g(\rho_i;A,c,\xi),
\]
where
\[
 g(x;A,c,\xi)=(A/\xi+1-c)x
 +\left[A+\frac\xi2+\frac{A^2-2c}{2\xi}\right]x^2
 -(c+Ac/\xi)x^3+\frac{c^2}{2\xi}x^4.
\]
First suppose $73/100\le k=\kappa\le7/9$ and let
$p^2=\sum_{\rho_i\ne k}\rho_i^4$ (the top eigenvalue is unique).
If $p\le3/20$, the matrix $B_0$ alone suffices. Indeed,
$s_3\le k^3+\sqrt{1-k^2}\,p$, $s_5\ge k^5$ and
$s_6\le k^6+(1-k^2)p^2$ give
\[
 \|B_0\|_F^2\le1+2k\sqrt{1-k^2}\,p-p^2,
\]
so its mixed cost is at most
\[
 k+kp^2+\sqrt{(k^4+p^2)(1+2k\sqrt{1-k^2}\,p-p^2)}
 <\frac{73}{50}.
\]
For the last comparison replace $k$ by $7/9$, $p$ by $3/20$,
and $2k\sqrt{1-k^2}$ by $1$, then square the positive sides.
The resulting rational inequality is verified by the checker.

If $p\ge3/20$, take
\[
 \begin{aligned}
 B&=\vartheta B_0+(1-\vartheta)
       [R+P((73/55)I-(103/55)R)],\\
 A&=\vartheta k+(1-\vartheta)73/55,\qquad
 c=\vartheta+(1-\vartheta)103/55.
 \end{aligned}
\]
The tangent condition is exact: $(73/55)^2=1+(48/55)^2$.
Write $d=k(1+Ak-ck^2)$ for the top eigenvalue of $B$, and
$q$ for the Frobenius norm of its remaining block.
Here $A^2\le c$, $d\ge k^2$ and
$q\le(5/4)\sqrt{1-k^2}<6/7$.
The function $x^{-1}+A-cx$ decreases, so
\[
 1\le\frac{k^2q}{dp}\le\frac{40}{7},\qquad
 \frac{k^2q}{dp}+\frac{dp}{k^2q}\le\frac{1649}{280}<6.
\]
Thus \eqref{eq:app-K-one-rankone-covariance} applies.
For any $\chi>0$, its mixed-cost bound and $2pq\le\chi p^2+q^2/\chi$ give
\[
 \mathbb E(TX^TBX)-\operatorname{tr}R
 \le A-ck^3+2k^2d+\frac{1-k^2}{\chi}
       +\sum_{\rho_i\ne k}\rho_i^2 n(\rho_i;A,c,\chi),
\]
where
\[
 n(x;A,c,\chi)=(2A/\chi-c)x+
       \left[\chi+\frac{A^2-2c}{\chi}\right]x^2
       -\frac{2Ac}{\chi}x^3+\frac{c^2}{\chi}x^4.
\]
Define the two scalar majorants
\[
 \begin{split}
 O(k,x)&=A+\frac1{2\xi}+k^2g(k;A,c,\xi)
                    +(1-k^2)g(x;A,c,\xi),\\
 N(k,x)&=A-ck^3+2k^2d+\frac{1-k^2}{\chi}
                    +(1-k^2)n(x;A,c,\chi).
 \end{split}
\]
Choose the parameters in the following table.
\[
\begin{array}{c|c|c|c|c}
 k\text{ interval}&\vartheta&\xi&\chi&\omega\\\hline
 {}[73/100,3/4]&9/20&8/5&2&3/10\\
 {}[3/4,7/9]&1/4&8/5&12/5&9/20
\end{array}
\]
For each row, the polynomial
$73/50-\omega O(k,x)-(1-\omega)N(k,x)$ has bidegree at most $(6,4)$.
All its tensor-product Bernstein coefficients exceed $1/5000$ on each
of the four rectangles obtained by taking
$x\in[7j/40,7(j+1)/40]$, $j=0,1,2,3$. Thus $280$ exact rational
comparisons certify both rows. Every remaining eigenvalue is at most
$\sqrt{1-k^2}<7/10$; averaging with weights
$\rho_i^2/(1-k^2)$ proves the desired mixed-cost bound.

For $k\ge7/9$, take $A=5/4$, $c=7/4$ and $\xi=8/5$.
Then $A+1/(2\xi)=25/16$ and
\[
 g(x)=\frac{x}{32}+\frac{1849x^2}{1280}
                 -\frac{399x^3}{128}+\frac{245x^4}{256}.
\]
All Bernstein coefficients of $7/100-g$ are positive on
$[0,1/4]$, $[1/4,3/8]$, $[3/8,1/2]$ and $[1/2,1]$.
If $J(k)=k^2g(k)+(1-k^2)7/100$, the Bernstein coefficients of
$-J'(k)/k$ on $[3/4,1]$ are also positive, so $J$ decreases there.
Consequently the cost is at most
\[
 \frac{25}{16}+J(7/9)
 <\frac{73}{50}.
\]
This completes the uniform-cost certificate.


\subsubsection{The lower correction budget}
Restore the scale $\eta$ and set
\[
 h(k)=\min\{2k,73/50\}.
\]
Use $B_0$ for $\kappa\le73/100$ and the uniform-cost comparison
for $\kappa\ge73/100$. They give the mixed-cost bound under the cap
\[
 \operatorname{tr}R+\eta h(\kappa)\le b.
\]
Lemma~\ref{lem:app-mixed-cap-cost} applies with
$a=1$, $f=0$, $o=2$ and $k_0=73/100$.
Put $k=73/100$. Its singular-value correction is bounded by
\[
 2\mathcal A(k)\le2\left(k-\frac{k^3}{6}-\frac{k^5}{40}\right)
 <\frac{33}{25}.
\]
Here we integrated $\sqrt{1-y^2}\le1-y^2/2-y^4/8$.
The latter inequality follows by squaring the positive right side:
$(1-z/2-z^2/8)^2-(1-z)=z^3/8+z^4/64\ge0$ for $0\le z\le1$.
All the polynomial certificates and the final rational rounding are checked by the
accompanying exact arithmetic script.



\subsubsection{The lower fluctuation bound}
Under the same cap, $\mathbb ET^2\le3$ allows us to take $m_2=3$ in
\eqref{eq:app-lower-moment-laplace}, giving
\begin{equation}\label{eq:app-K-one-laplace}
 \mathbb Ee^{-\lambda(1+b)\Delta_-}
 \le\exp\left\{-\frac{\lambda}{1+(5/2)\lambda}\right\}.
\end{equation}
Thus Lemma~\ref{lem:app-transfer} gives
$\beta=\sqrt{5/4}<28/25$ under the same cap, while the correction
budget gives $\alpha=33/25$.
These are the rounded $K=1$ lower-tail constants in
Theorem~\ref{thm:explicit-constants}.

% END K1 LOWER REFINEMENT

% BEGIN UPPER CONSTANTS
\subsection{The upper-tail constants}\label{app:upper-constants}

The exact upper shift is the soft upper spectral edge of $XX^T-R^{-1}$
at sensitivity $\operatorname{tr}R$, in the terminology of
Srivastava and Vershynin~\cite[Section~1.9, equation~(1.11)]{srivastava2013covariance}.
We prove the mean and exponential-moment bounds needed by
Lemma~\ref{lem:app-transfer} directly for the following parameters:
\[
\begin{array}{c|cc}
 &w&\nu_+\\\hline
 K\ge1&47/20&(9/2)K^2L\\
 K=1&58/25&4
\end{array}
\qquad L=\log(eK).
\]
Here $w$ is a dimensionless coefficient: the corresponding coefficient
in the cap of Lemma~\ref{lem:app-transfer} is $wK^2L$.
Throughout this subsection $R\succ0$, $\eta=\|R\|_F$, and
$K^2\operatorname{tr}R+wK^2L\eta\le b<1$.
Let $\Delta_+\ge0$ be the unique shift satisfying
\[
 R^{-1}+\Delta_+I-XX^T\succ0,\qquad
 \operatorname{tr}(R^{-1}+\Delta_+I-XX^T)^{-1}=\operatorname{tr}R.
\]
Our assertions are $\mathbb EY\le1$ for $Y=(1-b)\Delta_+$ and
\begin{equation}\label{eq:app-up-mgf}
 \mathbb Ee^{\lambda Y}\le
 \exp\!\left(\frac{\lambda}{1-\nu_+\lambda}\right),
 \qquad0<\lambda<1/\nu_+.
\end{equation}
The common spectral estimates come first, followed by the mean and
exponential-moment arguments. The mean bound controls $\gamma$, while
the MGF bound controls $\delta$; both use the same regularity cap.

\subsubsection{Exact shifts and their tails}
For $h\ge0$, put
\[
 R_h=(R^{-1}+hI)^{-1},\qquad
 B_h=hR_h+\frac{R_h^2}{\operatorname{tr}(RR_h)}.
\]
The Sherman--Morrison calculation in \papereqref{eq:shift-event} gives
\begin{equation}\label{eq:app-up-event}
 \{\Delta_+>h\}=\{X^TB_hX>h\},\qquad
 X^TB_{\Delta_+}X=\Delta_+.
\end{equation}

Write $r_j$ for the eigenvalues of $R$, and set
\[
 p_j(h)=\frac{r_j^2/(1+hr_j)}{\operatorname{tr}(RR_h)}.
\]
These numbers sum to one, and the corresponding eigenvalue of $B_h$ is
$(hr_j+p_j(h))/(1+hr_j)$.
For $0<q<1$,
\[
 1-q\frac{hr_j+p_j(h)}{1+hr_j}
 =\frac{1+(1-q)hr_j}{1+hr_j}
   \left(1-\frac{qp_j(h)}{1+(1-q)hr_j}\right).
\]
The inequalities
$\log(1+x)-\log(1+(1-q)x)\le qx$ for $x\ge0$ and
$\prod_j(1-a_j)\ge1-\sum_ja_j$ for $a_j\ge0$ with
$\sum_ja_j\le1$ therefore give
\[
 \det(I-qB_h)\ge(1-q)e^{-qh\operatorname{tr}R}.
\]
All the eigenvalue factors are positive. Apply the Gaussian comparison
of Lemma~\ref{app:gaussian-comparison} at $\lambda=q/(2K^2)$ and then
Markov's inequality in \eqref{eq:app-up-event}. Since
$1-K^2\operatorname{tr}R\ge1-b>0$, we obtain
\begin{equation}\label{eq:app-up-tail-family}
 \mathbb P\{Y>u\}\le(1-q)^{-1/2}e^{-qu/(2K^2)},
 \qquad u\ge0,\quad 0<q<1.
\end{equation}


We also use a more precise tail when the eigenvalues of $R$ are spread out.
Put $\eta=\|R\|_F$ and $\kappa=\|R\|/\eta\in(0,1]$.
The largest numerator $r_j^2/(1+hr_j)$ corresponds to $r_j=\|R\|$.
Since $\operatorname{tr}(RR_h)\ge\eta^2/(1+h\|R\|)$,
we have $p_j(h)\le\kappa^2$ for every $j$.
The concavity bound
$\log(1-qv)\ge(v/\kappa^2)\log(1-q\kappa^2)$ for
$0\le v\le\kappa^2$ improves the determinant estimate to
\[
 \det(I-qB_h)\ge(1-q\kappa^2)^{1/\kappa^2}
                          e^{-qh\operatorname{tr}R}.
\]
This holds also at $q=1$ when $\kappa<1$.
The same Gaussian comparison gives
\begin{equation}\label{eq:app-up-tail-kappa}
 \mathbb P\{Y>u\}\le
 (1-q\kappa^2)^{-1/(2\kappa^2)}e^{-qu/(2K^2)}.
\end{equation}
The margin in the cap allows a stronger version of this estimate:
\begin{equation}\label{eq:app-up-tail-margin}
 \mathbb P\{Y>u\}\le
 (1-q_0\kappa^2)^{-1/(2\kappa^2)}e^{-q_0u/(2K^2)}
\end{equation}
whenever $0<q_0\le2$, $q_0\kappa^2<1$, and
$u\kappa(q_0-1)\le wK^2L$.
To prove it, put
\[
 \vartheta=\frac{1-K^2\operatorname{tr}R}{1-b}
 \ge1+\frac{wK^2L\eta}{1-b},\qquad q=q_0/\vartheta,
 \qquad h=u/(1-b).
\]
If $q\le1$, the preceding determinant argument applies directly,
with a prefactor no larger than the one in \eqref{eq:app-up-tail-margin}.
If $q>1$, then $1\le\vartheta<q_0\le2$ and
\[
 h\|R\|\le\frac{\vartheta-1}{q_0-1}
 \le\frac{2-q}{q-1},\qquad
 \frac{q}{1-(q-1)h\|R\|}\le q_0.
\]
For the first comparison, after multiplying positive denominators the
difference is
$(q_0-1)(2-q_0)+(\vartheta-1)(\vartheta+q_0-2)\ge0$.
For the second, use
$(q-1)h\|R\|\le(\vartheta-1)/\vartheta$.
Now
$\log(1+x)-\log(1-(q-1)x)\le qx$ for
$0\le x\le(2-q)/(q-1)$, as differentiation shows.
Thus the first eigenvalue factors in the determinant still contribute
at least $e^{-qh\operatorname{tr}R}$.
The second factors have the form
$1-qp_j(h)/(1-(q-1)hr_j)$ and are at least $1-q_0p_j(h)$.
Their product is therefore at least
$(1-q_0\kappa^2)^{1/\kappa^2}$.
Gaussian comparison and Markov's inequality give exponent
$-q(1-K^2\operatorname{tr}R)h/(2K^2)=-q_0u/(2K^2)$,
which proves \eqref{eq:app-up-tail-margin}.

Put $\rho_j=r_j/\eta$ and $a=\sum_j\rho_j^3$, so that
$\kappa^3\le a\le\kappa$. For $q>0$ and $qk^2<1$, define
\begin{equation}\label{eq:k1-joint-prefactor}
 \mathcal G(k,a,q)=\exp\!\left\{
 \frac q2+\frac{a}{2k^3}
       [-\log(1-qk^2)-qk^2]\right\}.
\end{equation}
For $0<q\le2$, the tail comparison in \eqref{eq:app-up-tail-margin} improves to
\begin{equation}\label{eq:k1-joint-tail}
 \PP\{Y>u\}\le\mathcal G(\kappa,a,q)e^{-qu/(2K^2)},
 \qquad u\kappa(q-1)\le wK^2L.
\end{equation}
The sharper prefactor follows from majorization.
Order the $\rho_j$ decreasingly. The weights
\[
 p_j(h)=\frac{\rho_j^2/(1+h\eta\rho_j)}
                   {\sum_\ell\rho_\ell^2/(1+h\eta\rho_\ell)}
\]
are in the same order. Since the multiplier $(1+h\eta\rho_j)^{-1}$
increases with $j$, its weighted average on any initial segment is
at most its weighted average on the complementary segment.
Consequently every initial partial sum of $(p_j(h))$ is at most the
corresponding sum of $(\rho_j^2)$: the former vector is majorized by
the latter. Convexity of $v\mapsto v^r$, for integers $r\ge2$, gives
\[
 \sum_jp_j(h)^r\le\sum_j\rho_j^{2r}
 \le\kappa^{2r-4}\sum_j\rho_j^4
 \le a\kappa^{2r-3}.
\]
Expand $-\log(1-qp_j(h))$ into its positive power series and use
$\sum_jp_j(h)=1$. The resulting bound on the inverse square root of
$\prod_j(1-qp_j(h))$ is exactly \eqref{eq:k1-joint-prefactor}.
The proof of \eqref{eq:app-up-tail-margin} only uses this product after
bounding its second eigenvalue factors. Replacing that product bound
proves \eqref{eq:k1-joint-tail}; its margin calculation is unchanged
for every $K\ge1$.

To extend the tail estimate beyond rate two, define the scalar excess
\[
 F_v(\theta)=\log\frac{\theta(\theta+v-2)}{(\theta-1)^2+v-1}
            -\frac{v(\theta-1)}{\theta(v-1)}.
\]
The next lemma shows how an upper bound for this one-variable function
controls the loss in the determinant estimate. Its finite certificates
follow the proof.

\begin{lemma}[Upper-tail rates exceeding two]
\label{app:upper-tail-extension}
Suppose $K^2\operatorname{tr}R+wK^2L\eta\le b<1$ and
$0<\kappa=\|R\|/\eta\le k$.
Let $v>2$, $vk^2<1$, $C_v\ge0$, and suppose the scalar bound
$F_v(\theta)\le C_v$ holds for $1\le\theta\le v$.
Then, for every $1<v'\le v$,
\begin{equation}\label{eq:app-upper-extended-tail}
 \mathbb P\{(1-b)\Delta_+>u\}
 \le (1-vk^2)^{-1/(2k^2)}
       \exp\!\left(\frac{C_v}{2k^2}-\frac{v'u}{2K^2}\right)
 \quad\text{if }uk(v'-1)\le wK^2L.
\end{equation}
\end{lemma}

\begin{proof}
Set $Y=(1-b)\Delta_+$ and use the matrices $B_h$ and weights $p_j(h)$
from \eqref{eq:app-up-event}. Put
\[
 \vartheta=\frac{1-K^2\operatorname{tr}R}{1-b},\qquad
 h=\frac{u}{1-b},\qquad q=\frac{v'}{\vartheta}.
\]
The cap implies
$\vartheta-1\ge wK^2L\eta/(1-b)$.
If $q\le1$, the existing determinant argument applies immediately,
with a smaller prefactor than the one claimed.
Otherwise $1<\vartheta<v'$, and, with
$x_0=(\vartheta-1)/(v'-1)$, the assumed margin gives
\[
 h\eta\le x_0/k,\qquad hr_j\le(\rho_j/k)x_0\le x_0,
 \qquad \rho_j=r_j/\eta.
\]
All these numbers lie below $1/(q-1)$, because
\[
 (q-1)x_0
 =\frac{(v'-\vartheta)(\vartheta-1)}{\vartheta(v'-1)}
 \le\frac{\vartheta-1}{\vartheta}<1.
\]
In particular,
$q/[1-(q-1)hr_j]\le v'$.
The second factors in the determinant factorization are therefore
at least $1-v'p_j(h)>0$.

For the first factors define
\[
 F(q,x)=\log(1+x)-\log(1-(q-1)x)-qx.
\]
One has $F(q,0)=F_x(q,0)=0$ and
\[
 F_{xxx}(q,x)=\frac2{(1+x)^3}
          +\frac{2(q-1)^3}{(1-(q-1)x)^3}>0.
\]
The identity
$F(q,x)/x^2=\int_0^1(1-\omega)F_{xx}(q,\omega x)\,d\omega$
shows that $F(q,x)/x^2$ increases with $x$.
If $F(q,x_0)\ge0$, it follows that
\[
 \sum_jF(q,hr_j)
 \le\sum_j\frac{\rho_j^2}{k^2}F(q,x_0)
 =\frac{F(q,x_0)}{k^2}.
\]
If $F(q,x_0)<0$, every summand is nonpositive, so the sum is at most zero.
Thus in both cases the sum is at most
$\max\{0,F(q,x_0)\}/k^2$.

A direct substitution gives
\[
 F(q,x_0)=F_{v'}(\vartheta).
\]
To compare different rates, write $r=v-1$ and
$\upsilon=(\theta-1)/(v-1)\in[0,1]$. The definition of $F_v$ becomes
\[
 \log(1+\upsilon)+\log(1+r\upsilon)-\log(1+r\upsilon^2)
                    -\frac{(r+1)\upsilon}{1+r\upsilon},
\]
whose derivative with respect to $r$ equals
$r\upsilon^2(1-\upsilon)^2/[(1+r\upsilon)^2(1+r\upsilon^2)]\ge0$.
Therefore $F_{v'}(\vartheta)\le C_v$ whenever $1<v'\le v$.

Using $\sum_jp_j(h)=1$ and $p_j(h)\le\kappa^2$, the second determinant
factors have product at least $(1-v'\kappa^2)^{1/\kappa^2}$.
The logarithm of their inverse square root increases with both $v'$ and
$\kappa$, as follows by its positive power series, so the prefactor
is at most $(1-vk^2)^{-1/(2k^2)}$.
The first factors contribute at most
$qh\operatorname{tr}R/2+C_v/(2k^2)$ to the logarithm of the Gaussian
MGF bound. Markov's inequality now gives the claimed exponent, since
\[
 -\frac{q(1-K^2\operatorname{tr}R)h}{2K^2}
 =-\frac{v'u}{2K^2}.
\]

For the two-piece entropy tail, take
\[
 s=\frac{w}{k(v-1)},\qquad
 A=(1-vk^2)^{-1/(2k^2)}\exp\{C_v/(2k^2)\}.
\]
The choice $v'=v$ gives $Ae^{-vu/(2K^2)}$ up to $u=sK^2L$.
Above it choose $v'=1+s(v-1)K^2L/u\in(1,v]$.
The condition in \eqref{eq:app-upper-extended-tail} holds with equality,
and the bound becomes
$Ae^{-u/(2K^2)-s(v-1)L/2}$.
These two exponential pieces will be integrated in the entropy argument below.
\end{proof}

\paragraph{Certificates for the extended rates.}
The following choices satisfy $F_v(\theta)\le C_v$ on $[1,v]$:
\[
 (v,C_v)=(81/40,61/50000000),\quad
 (83/40,57/2000000),\quad(17/8,29/250000),
\]
\[
\begin{array}{c|rrrrrr}
 v&41/20&21/10&43/20&11/5&23/10&12/5\\\hline
 C_v&1/100000&13/200000&19/100000&1/2500&27/25000&21/10000\\[1mm]
 v&9/4&5/2&11/4&3&7/2&4\\\hline
 C_v&7/10000&7/2000&79/10000&27/2000&131/5000&201/5000
\end{array}
\]
For a finite rational certificate, set $y=\theta-1$, $V=vy$,
$W=2y^2+vy+2v-2$ and $A_1=(v-1)(1+y)$.
Then $0\le V/W<1$ and
$F_v(1+y)=\log[(1+V/W)/(1-V/W)]-vy/A_1$.
Use
\[
 U(z)=S_6(z)+\frac{2z^{13}}{13(1-z^2)}
 \ge\log\frac{1+z}{1-z},
\]
by \eqref{eq:app-arithmetic-log}. After multiplication by the positive denominator
$A_1W^{11}(W^2-V^2)$, the lower bound
$C_v+vy/A_1-U(V/W)$ has numerator
\[
 \begin{split}
 P_v(y)={}&(C_vA_1+vy)W^{11}(W^2-V^2)\\
 &-2A_1\left[
 \sum_{j=0}^{5}\frac{V^{2j+1}W^{10-2j}(W^2-V^2)}{2j+1}
             +\frac{V^{13}}{13}\right].
 \end{split}
\]
This is a degree-$27$ polynomial. Starting with $[0,v-1]$, test its
Bernstein coefficients and bisect every interval whose coefficients
are not all positive. At most eight bisections on any branch suffice
for the fifteen displayed rates; every terminal interval has positive
coefficients. The accompanying checker implements exactly this finite
rational certificate.

\subsubsection{The comparison matrix and determinant bounds}
The tail estimates alone do not give a sufficiently small mean cost.
We next bound the shift by a fixed quadratic comparison. Its determinant
excess $J$ measures the cost that will be combined with entropy to prove
$\mathbb EY\le1$.

\begin{lemma}[Fixed upper comparison]\label{app:upper-logdet}
Put $P=R^2/\eta^2$, $\bar r=\operatorname{tr}(RP)$, and
$C=R+P(\bar rI-2R)$. Then $C\succeq0$, $\|C\|<\eta/2$, and
\begin{equation}\label{eq:app-up-fixed-comparison}
 B_h\preceq P+hC\quad(h\ge0),\qquad
 \Delta_+\le X^TPX+\Delta_+X^TCX.
\end{equation}
For $a=\bar r/\eta$ the determinant excess
$J=-\log\det(I-2C/\eta)-2\operatorname{tr}R/\eta$ satisfies
\[
 J\le\mathcal J_0(a):=-2a+\frac{671}{100}
                    [1-a^2-a^2(1-a)^2].
\]
\end{lemma}
\begin{proof}
The earlier fixed comparison uses $R_h^2\preceq RR_h$, as in
\cite[Section~7.4, equation~(7.3)]{youssef2013estimating}.
Here we improve it using concavity of the individual eigenvalues of $B_h$.
Write $b_j(h)$ for the eigenvalue corresponding to $r_j$.
If $n=1$, then $B_h=P=1$ and $C=0$, so the assertions are immediate.
Otherwise put
\[
 S_j(h)=\sum_{k\ne j}\frac{r_k^2}{1+hr_k}.
\]
The eigenvalue formula gives
\[
 1-b_j(h)=\frac1{1+hr_j+r_j^2/S_j(h)}.
\]
Cauchy--Schwarz yields $2S_j'(h)^2\le S_j(h)S_j''(h)$,
so $1/S_j$ is concave. The reciprocal of a positive concave function is
convex; hence $b_j$ is concave. Its tangent at zero gives
\[
 b_j(h)\le b_j(0)+hb_j'(0)
 =p_j+h\bigl(r_j+p_j(\bar r-2r_j)\bigr),
 \qquad p_j=\frac{r_j^2}{\operatorname{tr}(R^2)}.
\]
This proves the matrix comparison, and substituting $h=\Delta_+$ gives the shift bound. Since $\bar r\ge p_jr_j$,
the eigenvalues of $C$ are at least $r_j(1-p_j)^2\ge0$.
Also $\operatorname{tr}C=\operatorname{tr}R-\bar r$.

Normalize $\eta=1$, write $x_j=r_j/\eta$, and put
$a=\sum_jx_j^3=\bar r/\eta\in(0,1]$.
The normalized eigenvalues of $C$ are
$c_j=x_j+ax_j^2-2x_j^3$.
We first record two norm bounds:
\begin{equation}\label{eq:app-up-comparison-norms}
 \sum_jc_j^2\le1-a^2-a^2(1-a)^2,
 \qquad \max_jc_j<\frac{223}{500}.
\end{equation}
For the first, put $p_j=x_j^2$ and define the symmetric nonnegative matrix
\[
 A_{jj}=(1-p_j)^2,\qquad A_{jk}=p_jp_k\quad(j\ne k).
\]
Then $c=Ax$ and the row and column sums of $A$ are $1-p_j$.
Jensen's inequality after normalizing each nonzero row gives
\[
 \sum_jc_j^2
 \le\sum_j(1-p_j)\sum_k A_{jk}x_k^2
 =1-2s_2-s_2^2+2s_3,
 \quad s_2=\sum_jp_j^2,\quad s_3=\sum_jp_j^3.
\]
Now $s_3\le s_2^{3/2}$ and $s_2\ge a^2$.
The function $u\mapsto1-2u-u^2+2u^{3/2}$ decreases on $[0,1]$,
which proves the first bound in \eqref{eq:app-up-comparison-norms}.
A row of sum zero contributes zero and requires no division.

For the second bound, fix $x=x_j$ and let $y=\sqrt{1-x^2}$.
Since $a\le x^3+y^3$,
\[
 c_j\le x+x^2(x^3+y^3)-2x^3=xy^3(x+y).
\]
Writing $u=x/y$ when $y>0$, this is
$u(1+u)/(1+u^2)^{5/2}$.
Its derivative has the sign of $-(3u+1)(u^2+u-1)$, so its maximum is
$((5+\sqrt5)/10)^{5/2}$.
The square of this maximum is
\[
 \frac{25+11\sqrt5}{250}<\frac{1241}{6250}
 <\left(\frac{223}{500}\right)^2.
\]
Here $\sqrt5<56/25$ suffices. The case $y=0$ gives $C=0$.
This proves \eqref{eq:app-up-comparison-norms}, and in particular
$I-2C\succ0$.

\emph{An analytic bound when $a$ is large.}
Put $M=223/500$. Expanding the logarithm and using
$\sum c_j=\sum x_j-a$ and $c_j\le M$ gives
\[
 J\le-2a+
 \frac{-\log(1-2M)-2M}{M^2}\bigl[1-a^2-a^2(1-a)^2\bigr].
\]
The displayed coefficient is less than $671/100$.
Indeed, $1-2M=27/250$ and the positive Taylor sum for
$e^{222672636/10^8}$ through degree eight exceeds $250/27$, while
$222672636/10^8=2M+(671/100)M^2$.
\end{proof}

\paragraph{Six scalar majorants for the determinant excess.}
The preceding norm bound is useful for large $a$. On other spectral
ranges we obtain smaller bounds by summing scalar logarithmic majorants.
Each is used only on the domain certified below.

Write $x_j=r_j/\eta$ and $a=\sum_jx_j^3$.
All the following scalar formulas use the normalization $\sum_jx_j^2=1$.
Set $v_a=a-7/20$ and
\[
\begin{array}{c|c|c}
 i&\mathcal B_i(a)&\mathcal A_i(a)\\\hline
 1&13/6&1+9a\\
 2&13/6-4(a-1/4)^2&20a-2\\
 3&13/5&24a-22/3\\
 4&13/6-4(a-1/4)^2-1/30&20a-2\\
 5&13/6-1/40&1+9a
\end{array}
\]
\[
\begin{aligned}
 \mathcal B_6(a)&=\frac{413}{200}-\frac{1339}{1000}v_a
               -\frac{10009}{1000}v_a^2-\frac{25197}{1000}v_a^3,\\
 \mathcal A_6(a)&=\frac{816}{125}+\frac{7307}{200}v_a
                +\frac{9971}{100}v_a^2+\frac{18903}{250}v_a^3,\\
 \mathcal H_i(a,x)&=2x+\mathcal B_i(a)x^2-\mathcal A_i(a)x^2(x-a),\qquad
 z(a,x)=2(x+ax^2-2x^3).
\end{aligned}
\]
The six inequalities
\begin{equation}\label{eq:app-up-scalar-logdet}
 -\log(1-z(a,x))\le\mathcal H_i(a,x)
\end{equation}
hold on the domains in the table below. Their common rational certificate
uses
\[
\begin{aligned}
 \mathcal D(z)&=1-\frac{16}{7}z+\frac{12}{7}z^2
                    -\frac{16}{35}z^3+\frac1{35}z^4,\\
 \mathcal N(z)&=z-\frac{25}{14}z^2+\frac{19}{21}z^3
                    -\frac{47}{420}z^4.
\end{aligned}
\]
Indeed,
\[
 \mathcal D(z)=\frac{1-z}{35}
 [4+18(1-z)+12(1-z)^2+(1-z)^3]>0\qquad(z<1),
\]
and direct differentiation gives
\[
 \frac{d}{dz}\left[\frac{\mathcal N(z)}{\mathcal D(z)}+\log(1-z)\right]
 =\frac{z^7}{1225\mathcal D(z)^2}.
\]
The bracket vanishes at zero and is nonnegative on either side of zero.
Consequently $-\log(1-z)\le\mathcal N(z)/\mathcal D(z)$ for $z<1$.
On all domains below, $z<1$: for $x\le1/2$ the largest cubic
$x+(5/6)x^2-2x^3$ is at most $11/24$, and for $x=1/2+u$ it is at most
$11/24+u/3-2u^2\le17/36$.

The polynomials
$\mathcal Q_i=(\mathcal H_i\mathcal D(z)-\mathcal N(z))/x^2$
extend to $x=0$ by cancellation. The following table specifies their
initial rectangles, bidegrees, and the number of terminal rectangles
in a single uniform certificate:
\[
\begin{array}{c|c|c|c|r}
 i&a\text{ interval}&x\text{ interval}&\text{bidegree}&\text{rectangles}\\\hline
 1&[0,27/40]&[0,1]&(6,13)&19\\
 2&[1/4,37/50]&[0,1]&(6,13)&19\\
 3&[7/10,5/6]&[0,1]&(6,13)&13\\
 4&[1/4,7/10]&[0,1]&(6,13)&28\\
 5&[0,13/20]&[0,1]&(6,13)&16\\
 6&[1/4,51/100]&[0,13/20]&(8,13)&70
\end{array}
\]
Use the tensor-product Bernstein coefficients defined above.
Starting from each displayed domain, retain a
rectangle if all its coefficients are positive; otherwise bisect both
coordinates and test all four children. The exact checker reproduces
this deterministic rule. It stops within five bisections on every branch,
with the displayed counts, and every terminal coefficient exceeds
$1/100000$. This proves all six scalar majorants.

Summing $\mathcal H_i(a,x_j)$ cancels the term containing $x_j-a$.
Retaining the scalar excess at a largest eigenvalue gives
\begin{equation}\label{eq:k1-joint-J}
 J\le\mathcal B_i(a)-\mathcal H_i(a,\kappa)
       +\frac{\mathcal N(z(a,\kappa))}{\mathcal D(z(a,\kappa))}.
\end{equation}
Use a majorant only on its stated $a$ interval; the sixth one additionally
requires $\kappa\le13/20$. For $a\ge5/6$ the norm bound
$\mathcal J_0(a)$ supplies the alternative certificate.
In particular, summing the fifth majorant gives $J\le257/120$ for
$a\le13/20$.

We first record an additional determinant estimate for concentrated spectra.
If $\kappa\ge37/40$, every eigenvalue $\rho_i$ except a largest one is
less than $381/1000$, since
$1-\kappa^2\le231/1600<(381/1000)^2$.
As $a\le1$, its contribution to $C/\eta$ is at most
$f(\rho_i)$, where $f(x)=x+x^2-2x^3$.
This function increases on $[0,381/1000]$ and
$f(381/1000)<21/50$.
The largest contribution is at most
$\kappa+\kappa^2-2\kappa^3\le3(1-\kappa)\le9/40$.
Thus $\|C\|/\eta\le M:=21/50$.
The Frobenius-norm argument in Lemma~\ref{app:upper-logdet} yields
\begin{equation}\label{eq:general-high-kappa-J}
 J\le\mathcal J_1(a):=-2a+\frac{563}{100}
                 [1-a^2-a^2(1-a)^2]
 \qquad(\kappa\ge37/40).
\end{equation}
Indeed, the positive degree-$24$ Taylor sum proves
\[
 e^{2M+(563/100)M^2}>25/4=1/(1-2M),
\]
which gives the claimed logarithmic coefficient.


\subsubsection{Entropy bounds and the upper mean}
We now combine the determinant estimates with the shift tails. The common
reduction below shows that it suffices to bound an entropy cost plus $J/2$;
the two mean lemmas then certify this bound at $K=1$ and for general $K$.

For any mean-one $W\ge0$, with $0\log0=0$, write $D(W)=\mathbb EW\log W$.
For $c>0$, the part below $c$ of $W\log(W/c)$ is nonpositive, so
\begin{equation}\label{eq:app-up-entropy-integration}
 D(W)\le\log c+\int_c^\infty[1+\log(u/c)]\mathbb P\{W>u\}\,du.
\end{equation}
For later integration put $I(x)=\int_0^\infty e^{-u}/(x+u)\,du$, $x>0$.
The entropy integral is bounded using
\[
 \frac{x^2+8x+11}{x^3+9x^2+18x+6}
 \le I(x)\le
 \frac{x^3+11x^2+26x+6}{x(x^3+12x^2+36x+24)}.
\]
For completeness, with $U$ exponential of mean one,
$D_0=x^3+9x^2+18x+6$ and $D_1=x^3+12x^2+36x+24$, the exact remainders are
\[
 \begin{aligned}
 I(x)-\frac{x^2+8x+11}{D_0}
 &=\frac1{D_0^2}\mathbb E
       \frac{(U^3-9U^2+18U-6)^2}{x+U},\\
 \frac{x^3+11x^2+26x+6}{xD_1}-I(x)
 &=\frac1{xD_1^2}\mathbb E
       \frac{U(U^3-12U^2+36U-24)^2}{x+U}.
 \end{aligned}
\]
Expanding these identities uses only $\mathbb EU^j=j!$;
the accompanying checker verifies the polynomial and factorial-moment
identities exactly.

For both upper-mean certificates, use
\eqref{eq:app-arithmetic-exp}--\eqref{eq:app-arithmetic-log}
with $18$ logarithmic terms after reduction by powers of two, and
exponential Taylor degree $80$.
Round intermediate bounds outwards to multiples of $10^{-15}$.
On the $K=1$ integration grid, successive $e^{-qu_j/2}$ values use
multiplication by an upper bound for $e^{-q/400}$ and upward rounding.
The final rounding precisions are specified below.

The same final step converts both entropy estimates below into mean bounds.
If $\mu=\mathbb EY\ge1$, then $W=Y/\mu\le Y$ has mean one and inherits
all its tail bounds. If $D(W)+J/2\le wL$, Gaussian comparison and
\eqref{eq:app-entropy-variational} give
\[
 \mathbb E(WX^TCX)\le K^2\eta\bigl(D(W)-\tfrac12\log\det(I-2C/\eta)\bigr)
 \le K^2\operatorname{tr}R+wK^2L\eta\le b.
\]
The fixed comparison then gives $\mu\le1-b+b\mu$, hence $\mu\le1$.
The tail family ensures that $\mu$ is finite.

\begin{lemma}[Upper mean when $K=1$]\label{lem:k1-upper-entropy}
For $K=1$ the cap $\operatorname{tr}R+(58/25)\eta\le b<1$ implies
$\mathbb EY\le1$.
\end{lemma}
\begin{proof}
Set $w=58/25$. By the preceding reduction it suffices to prove
\begin{equation}\label{eq:k1-joint-cost}
 D(W)+J/2\le w
\end{equation}
for every mean-one weight satisfying \eqref{eq:k1-joint-tail}.
First suppose $0<\kappa\le1/4$. The prefactor in
\eqref{eq:k1-joint-prefactor} is at most
$(1-q\kappa^2)^{-1/(2\kappa^2)}$, because $a\le\kappa$.
On each of the fifty intervals $[i/200,(i+1)/200]$, write
$k_i=(i+1)/200$ and choose $c_i$ by rounding
$43/20+(23/10)k_i^2$ upwards to a multiple of $1/200$.
Use the parameters $q\in\{j/20:18\le j\le40\}$ satisfying
$qk_i^2<1$ and $u_jk_i(q-1)\le w$, with $u_j=j/200$.
The upper Riemann sum for \eqref{eq:app-up-entropy-integration} uses the
minimum of these tail bounds and $1$ at $u_j$, multiplied by
$[1+\log(u_{j+1}/c_i)]/200$.
Integrate this sum to $40$. For the remaining integral use
$q=9/10$, $\sqrt{10}<31623/10000$ and $\log(u/c_i)\le u/c_i$.
Writing $\omega=9/20$, the remainder is at most
\begin{equation}\label{eq:k1-entropy-remainder}
 \mathcal R_0(c)=\frac{31623}{10000}e^{-40\omega}
 \left(\frac1\omega+\frac{40}{c\omega}+\frac1{c\omega^2}\right).
\end{equation}
After adding $\log c_i$, each resulting bound lies more than $1/500$ below
\[
 \frac{239}{200}+\frac9{20}\left(\frac i{200}\right)^2
       +\frac12\left(\frac i{200}\right)^4.
\]
The polynomial is increasing. Together with $J\le13/6-1/40$, proved above, this gives
\[
 D(W)+\tfrac12J\le
 \frac{239}{200}+\frac{9}{20}\left(\frac14\right)^2
       +\frac12\left(\frac14\right)^4+\frac{13}{12}-\frac1{80}<\frac{58}{25}.
\]

For $\kappa\ge1/4$, use the following joint certificate.
Partition $[1/4,1]$ at $k=i/80$ for $20\le i\le80$.
Partition $[0,1]$ at all $j/40$, adding $37/50$ and $5/6$.
Discard rectangles $[a_-,a_+]\times[k_-,k_+]$ for which
$a_+<k_-^3$ or $a_->k_+$.
Increase every remaining lower $a$ endpoint to
$\max\{a_-,k_-^3\}$ and discard empty intervals.
The resulting $837$ rectangles still cover all admissible pairs.
The joint polynomial test below succeeds on all $837$ rectangles
without any further subdivision.

For each rectangle, let $\bar k,\bar a$ denote its midpoints and put
$v=\bar k-1/2$. Set
\[
 c_0=\begin{cases}
 43/20+(6/5)\bar k^2+(9/2)\bar k^4,&\bar k\le1/2,\\
 273/100+(37/10)v+17v^2-11v^3,&\bar k>1/2.
 \end{cases}
\]
Round $c_0-(2\bar k+1/4)(\bar k-\bar a)$ upwards to the next
multiple of $1/200$ to obtain $c$.
For $u_j=j/200$, $200c\le j<8000$, put
\[
 \begin{aligned}
 q_j^*&=\frac{u_j-1}{\bar k\bar a+\bar k^2(u_j-1)},\\
 q_j&=\frac1{100}\max\left\{90,
 \min\left\{200,\lfloor100q_j^*\rfloor,
 100+\left\lfloor\frac{100w}{u_jk_+}\right\rfloor,
 \left\lceil\frac{100}{k_+^2}\right\rceil-1\right\}\right\}.
 \end{aligned}
\]
These parameters satisfy all restrictions in
\eqref{eq:k1-joint-tail} at $u_j$ throughout the rectangle.
Define
\[
 E(a,k)=\log c+\frac1{200}\sum_{j=200c}^{7999}
 [1+\log(u_{j+1}/c)]\mathcal G(k,a,q_j)e^{-q_ju_j/2}
 +\mathcal R_0(c).
\]
The monotonicity of survival probabilities gives $D(W)\le E(a,\kappa)$.
The separate convexity of $E$ is useful here: indeed,
\[
 \log\mathcal G(k,a,q)=\frac q2+\frac a2
                \sum_{r=2}^{\infty}\frac{q^r k^{2r-3}}r
\]
is affine in $a$ and convex in $k$, with nonnegative coefficients.
Exponentiation preserves both separate convexities.
Thus $E$ lies below the bilinear interpolation of its four corner
values. Bound those values upwards, rounding finally to multiples of
$10^{-9}$, and denote their bilinear interpolation by $\mathcal E_0$.

On each rectangle, use any one of the applicable bounds
\eqref{eq:k1-joint-J}. The checker verifies that all Bernstein
coefficients of at least one applicable polynomial
\[
 \left[w-\mathcal E_0-\frac{\mathcal B_i}{2}
                      +\frac{\mathcal H_i}{2}\right]\mathcal D
                 -\frac{\mathcal N}{2}
\]
are strictly positive. On rectangles with $a_-\ge5/6$, it can instead
verify the polynomial $w-\mathcal E_0-\mathcal J_0/2$ with strictly positive coefficients. Since the relevant denominator $\mathcal D$ is
positive, these checks establish \eqref{eq:k1-joint-cost}.
The rational logarithm and exponential rules above evaluate all corner bounds;
both spectral ranges use the same $1/200$ integration mesh.
The common mean reduction completes the proof.
\end{proof}

\begin{lemma}[Upper mean for general $K$]\label{app:general-upper-drift}
Let $K\ge1$, $L=\log(eK)$ and $w=47/20$.
In the notation of the upper-shift proof, suppose
\[
 K^2\operatorname{tr}R+wK^2L\eta\le b<1,
 \qquad \eta=\|R\|_F.
\]
Then $Y=(1-b)\Delta_+$ has mean at most one.
\end{lemma}

\begin{proof}
Set $w=47/20$. On a spectral rectangle with upper endpoints $k,a_+$,
use the tail prefactor
\[
 A=\mathcal G(k,a_+,q)e^{C_q/(2k^2)},\qquad
 s=\frac{w}{k(q-1)}\quad(q>1),
\]
where $C_q=0$ for $q\le2$ and the rate certificate supplies it for $q>2$.
The positive series for $\log\mathcal G$ shows that it increases in each
argument, so this prefactor works also for the adaptive rate above the cutoff.
Integrating the resulting tails in \eqref{eq:app-up-entropy-integration},
with cutoff $cK^2L$, gives
\[
 D(W)\le2L-2+\log(cL)+\mathcal E(L),\qquad \mathcal E=\mathcal R+\mathcal T,
\]
where
\[
\begin{aligned}
 \mathcal R(L)&=\frac{2A}{q}e^{-2-(qc/2-2)L}[1+I(qcL/2)],\\
 \mathcal T(L)&=Ae^{-2-(qs/2-2)L}
 \left[(2-2/q)(1+\log(s/c))+2I(sL/2)-\frac2qI(qsL/2)\right].
\end{aligned}
\]
Set $\mathcal T=0$ for $q\le1$.
Every witness has $qc/2>2$ and, for $q>1$, also $s\ge c$ and $qs/2>2$.
Thus $\mathcal E$ is decreasing and convex. For $\mathcal R$ this follows
from the integral defining $I$; for $\mathcal T$ use
\[
 I(x)-\frac1qI(qx)=\int_0^\infty e^{-xv}
 \left[\frac1{1+v}-\frac1{q(q+v)}\right]\,dv\qquad(q>1).
\]
The kernel is nonnegative, and multiplication by a decreasing exponential
preserves these properties.

For $\kappa\le2/5$ we use the fifth scalar majorant
$\mathcal H_5$, which strengthens the first and is already proved for $a\le13/20$.
Summing it cancels the term containing $\rho_i-a$, and gives
$J\le13/6-1/40=257/120$.
Use $k=a_+=2/5$, $q=11/4$ and $c=193/100$ in this range.
The entropy calculation below has margin greater than $1/100000$.

For $2/5\le\kappa\le1$, start with the intervals
$[i/80,(i+1)/80]$, $32\le i\le79$.
Partition $a\in[0,1]$ at the multiples of $1/40$, adding
$27/40$, $37/50$ and $5/6$.
On each rectangle $[k_-,k_+]\times[a_-,a_+]$, impose the spectral
restrictions by replacing $k_-$ by $\max\{k_-,a_-\}$,
then $a_-$ by $\max\{a_-,k_-^3\}$ with this updated endpoint,
and $a_+$ by $\min\{a_+,k_+\}$.
Discard rectangles with empty interior; boundary cases follow by continuity.
Some rectangles are subdivided at both midpoints, with the same
restrictions imposed on each child. At most five such bisections on any
branch suffice.
The exact checker supplies and verifies rational witnesses on the resulting
$1330$ rectangles, and checks that these rectangles cover every initial
rectangle after the restrictions.

A witness consists of a rational determinant bound $J_*$, an applicable
scalar majorant, a rate $q$, and a cutoff coefficient $c$.
The permitted rates are the multiples of $1/40$ in $[9/10,2]$ together
with the fifteen rates in the extension certificate.
The cutoffs are positive multiples of $1/100$ not exceeding $8$.
For each of the six scalar majorants, the determinant certificate is
positivity of every Bernstein coefficient of
\[
 [J_*-\mathcal B_i(a)+\mathcal H_i(a,\kappa)]
       \mathcal D(z(a,\kappa))-\mathcal N(z(a,\kappa))
\]
on the rectangle, with the majorant used only on its established domain.
The coefficients of $\mathcal D$ are also checked positive after degree
elevation to the same bidegree.
Alternatively, when $a_-\ge5/6$ use the polynomial
$J_*-\mathcal J_0(a)$ from the earlier norm bound;
when $k_-\ge37/40$ one may use
$J_*-\mathcal J_1(a)$ from \eqref{eq:general-high-kappa-J}.
These finite checks prove $J\le J_*$ throughout every rectangle.
Use the upper bound for positive coefficients of $I$ and the lower bound
for its negative coefficient in $\mathcal T$.
At $L=j/40$, $40\le j\le160$, apply the rational logarithm and exponential
rules above, rounding intermediate bounds outwards
to multiples of $10^{-15}$ and final upper bounds for $\mathcal E(L)$
to multiples of $10^{-10}$.
Write $e_\ell,e_h$ for consecutive such bounds and
$\sigma=(e_h-e_\ell)/(h-\ell)$.
Convexity bounds $\mathcal E$ by the chord joining them.
For $r=w-2-\sigma>0$ put
$x_*=\min\{h,\max\{\ell,1/r\}\}$.
The finite comparison
\[
 2+rx_* -\log(cx_*)-\frac{J_*}{2}-e_\ell+\sigma\ell
 >\frac1{100000}
\]
then proves $D(W)+J/2<wL$ on $[\ell,h]$.
For $L\ge4$, use $\mathcal E(L)\le e_4$ and $1/(w-2)<4$;
the remaining comparison is
\[
 2+4(w-2)-\log(4c)-\frac{J_*}{2}-e_4>\frac1{100000}.
\]
The checker verifies these inequalities for every rational witness.
This proves the required entropy bound for all $L\ge1$ and all spectra.

The common mean reduction completes the proof.
\end{proof}

\subsubsection{Exponential moments}
With the mean bounds established, we turn to $\nu_+$ and the fluctuation
constant. A first-moment bound and an exponential tail control all higher
moments. The second moment needs a separate refinement, especially near
$K=1$ and for concentrated spectra.

\begin{lemma}[Moments from a first moment and an exponential tail]
\label{app:upper-scalar-mgf}
Suppose $Y\ge0$, $\mathbb EY\le1$, and
$\mathbb P\{Y>u\}\le\sqrt{10}e^{-9u/(20K^2)}$.
Then $\mathbb EY^p/p!\le[(9/2)K^2L]^{p-1}$ for every integer $p\ge3$.
For $L\ge5/4$ one also has
$\mathbb EY^2/(2K^2L)-1/(2K^2L)<9/2$.
\end{lemma}
\begin{proof}
We combine the first moment and the tail as in
\cite[Lemma~2.1]{jeong2022subgaussian}.
Put
$A_*=15K^2L/2$, $s_0=20K^2/9$, and
$\varepsilon_0=\sqrt{10}\,s_0e^{-A_*/s_0}$.
Since $L\ge1$,
\[
 \varepsilon_0=\frac{20\sqrt{10}}9e^{-2-11L/8}
 \le\frac14 e^{-11(L-1)/8}\le\frac14.
\]
For the first inequality, use $\sqrt{10}<3163/1000$ and
$e^{27/8}>146/5$; the latter follows from the positive Taylor sum
through degree fourteen.
For each integer $p\ge2$, the function
$u\mapsto u^p-A_*^{p-1}u$ is nonpositive below $A_*$.
Integrating its positive part against the tail and using $\mathbb EY\le1$ gives
\begin{equation}\label{eq:app-tail-moment-integration}
 \frac{\mathbb EY^p}{p!}
 \le(1-\varepsilon_0)\frac{A_*^{p-1}}{p!}
   +\varepsilon_0\sum_{j=0}^{p-1}\frac{A_*^j s_0^{p-1-j}}{j!}.
\end{equation}
The coefficient of $\varepsilon_0$ is positive. Since
$s_0\le(20/9)K^2L$, it follows that
\[
 \frac{\mathbb EY^p}{p!}\le(K^2L)^{p-1}u_p,
 \qquad
 u_p=\frac34\frac{(15/2)^{p-1}}{p!}
     +\frac14\sum_{j=0}^{p-1}
       \frac{(15/2)^j(20/9)^{p-1-j}}{j!}.
\]
Direct rational arithmetic gives
\[
 u_3=\frac{25225}{1296}<\left(\frac92\right)^2,
 \qquad
 u_4=\frac{10896875}{186624}<\left(\frac92\right)^3.
\]
For $p\ge4$, the first component grows by the factor $(15/2)/(p+1)$.
Writing $T_p$ for the sum in the second component,
\[
 T_{p+1}=\frac{20}{9}T_p+\frac{(15/2)^p}{p!}
 \le\left(\frac{20}{9}+\frac{15}{2p}\right)T_p
 \le\frac{295}{72}T_p<\frac92T_p.
\]
This proves the higher-moment assertion.

For the second moment the same integration gives
\[
 \frac{\mathbb EY^2}{2K^2L}-\frac1{2K^2L}
 \le f_0(L):=\frac{15}{4}+\frac14e^{-11(L-1)/8}
       \left(\frac{15}{4}+\frac{20}{9L}\right)
       -\frac{e^{-2(L-1)}}{2L}.
\]
For $L\ge1$,
\[
 -L^2e^{2(L-1)}f_0'(L)
 \ge\frac{165}{128}L^2-\frac{17}{72}L+\frac1{18}>0.
\]
The polynomial has positive coefficients in $L-1$.
The bounds \eqref{eq:app-arithmetic-exp} with $N=24$ at $11/32$ and
$1/2$ give $f_0(5/4)<9/2$.
These are rational checks in the accompanying script.
\end{proof}

\begin{lemma}[Upper exponential moment when $K=1$]\label{app:one-upper-mgf}
Assume $K=1$ and
$\operatorname{tr}R+w\eta\le b<1$, where $\eta=\|R\|_F$ and $w\ge23/10$.
Suppose the normalized upper shift $Y=(1-b)\Delta_+$ satisfies
$\mathbb EY\le1$. Then \eqref{eq:app-up-mgf} holds with $\nu_+=4$.
\end{lemma}

\begin{proof}
Put $\kappa=\|R\|/\eta$.
It suffices to bound the moment generating function by
$1+\lambda+\sum_{p\ge2}u_p\lambda^p$, where
$u_2\le9/2$ and $u_p\le4^{p-1}$ for every integer $p\ge3$.
Indeed, this series is at most
\[
 1+\frac{\lambda}{1-4\lambda}+\frac{\lambda^2}{2}
 \le\exp\left(\frac{\lambda}{1-4\lambda}\right)
 \qquad(0<\lambda<1/4).
\]

First suppose $\kappa\le15/16$ and put $v=225/256$.
The weights $p_j(h)$ in the determinant argument obey
$0\le p_j(h)\le v$ and $\sum_jp_j(h)=1$.
For any such weights, $\prod_j(1-p_j(h))\ge v(1-v)$.
Indeed, replacing two weights by $0,a+b$ when $a+b\le v$, or by
$a+b-v,v$ otherwise, decreases the product. Repetition leaves just
$v,1-v$ as the nonzero weights, since $v>1/2$.
As $1/[v(1-v)]<(46/15)^2$, the determinant argument at $q=1$ gives
$\mathbb P\{Y>u\}\le(46/15)e^{-u/2}$.
The first-moment and tail integration in
Lemma~\ref{app:upper-scalar-mgf}, with cutoff $A=13/2$ and scale $s=2$,
has tail integral $(46/15)s e^{-A/s}\le5/21$.
Thus
\[
 \frac{\mathbb EY^p}{p!}\le u_p
 :=\frac{16}{21}\frac{A^{p-1}}{p!}
 +\frac5{21}\sum_{j=0}^{p-1}\frac{A^j s^{p-1-j}}{j!}.
\]
Direct rational arithmetic gives $u_2=9/2$, $u_3<16$, and $u_4<64$.
For $p\ge4$, the two components grow by factors at most
$A/(p+1)<4$ and $s+A/p\le29/8<4$, respectively.
This proves all required bounds in this case.

For the remaining spectra, set
\[
 P=R^2/\eta^2,\quad T=X^TPX,\quad
 \bar r=\operatorname{tr}(RP),\quad
 C=R+P(\bar rI-2R),\quad Q=X^TCX.
\]
Lemma~\ref{app:upper-logdet} gives $C\succeq0$,
$\operatorname{tr}C=\operatorname{tr}R-\bar r$, and
$\Delta_+\le T+\Delta_+Q$.
Write $\rho_i$ for the eigenvalues of $R/\eta$ and $a=\sum_i\rho_i^3$;
the eigenvalues of $C/\eta$ are $c_i=\rho_i+a\rho_i^2-2\rho_i^3$.

We first record a scalar monotonicity fact. Put
$\mathcal J(z)=-\log(1-z)-z$ and $z=(12/5)(x+Ax^2-2x^3)$.
Then
\begin{equation}\label{eq:app-one-upper-logratio}
 x\longmapsto\frac{\mathcal J((12/5)(x+Ax^2-2x^3))}{x^2}
 \quad\hbox{is increasing on }(0,7/20],\qquad 1/2\le A\le1.
\end{equation}
Here $0<z<1$: the defining cubic is increasing in $x$ and $A$
on the stated rectangle, and its maximum gives $z<1$.
To verify the derivative sign using the existing rational
logarithm bound $-\log(1-z)\le\mathcal N(z)/\mathcal D(z)$, it suffices
that
\[
 xz'(x)z\mathcal D(z)
 -2\bigl(\mathcal N(z)-z\mathcal D(z)\bigr)(1-z)\ge0.
\]
After division by $x^3$, this is a polynomial in $(A,x)$.
Its Bernstein coefficients on $[1/2,1]\times[0,7/20]$ are all greater
than $49/500$, by the coefficient rule already given above.
This is a finite rational certificate of \eqref{eq:app-one-upper-logratio}.

Use the following bounds in the three indicated spectral intervals:
\[
\begin{array}{c|c|c|c|c|c}
 \text{interval for }\kappa&r&A&q_*&c_*&d\rule{0pt}{3ex}\\ \hline
 {}[15/16,47/50]&87/250&1087991/1250000&13/250&37/100&31239/10000\\
 {}[47/50,19/20]&171/500&111/125&1/20&183/500&313/100\\
 {}[19/20,1]&5/16&1&1/25&7/20&63/20
\end{array}
\]
In every row, the largest eigenvalue of $R/\eta$ contributes at most
$q_*$ to $C/\eta$; the other eigenvalues satisfy $\rho_i\le r$,
$\sum_{i\ne\mathrm{top}}\rho_i^2\le r^2$, and
$c_i\le\rho_i+A\rho_i^2-2\rho_i^3$.
Moreover $r+Ar^2-2r^3<c_*$ and
$b-\operatorname{tr}C\ge d\eta$.

For clarity, these bounds follow directly from the endpoint spectra.
On $[\kappa_{\rm lo},\kappa_{\rm hi}]$ in the first two rows put
$\varepsilon=1-\kappa_{\rm lo}^2$.
The displayed $r$ exceeds $\sqrt\varepsilon$.
Since $z^3+(1-z^2)^{3/2}$ increases for $z\ge1/\sqrt2$,
\[
 a\le\kappa_{\rm hi}^3+(1-\kappa_{\rm hi}^2)^{3/2}\le A.
\]
For the first row use $\sqrt{1-(47/50)^2}<171/500$; for the second
use $\sqrt{1-(19/20)^2}<5/16$.
The largest contribution is at most
$\kappa_{\rm hi}\varepsilon^2+\kappa_{\rm hi}^2\varepsilon r<q_*$.
For the last row it is at most
$\varepsilon^2+(5/16)\varepsilon<1/25$, where
$\varepsilon\le39/400$.
Finally $a\ge\kappa^3$ gives the gap from
$23/10+\kappa_{\rm lo}^3>d$ in each row.
All stated comparisons are rational.

Fix $\theta_0=6/5$.
Gaussian comparison and \eqref{eq:app-one-upper-logratio} now give
\[
 \begin{split}
 \log\mathbb E\exp\left(\frac{\theta_0(Q-\operatorname{tr}C)}\eta\right)
 &\le\frac12\sum_i\mathcal J(2\theta_0c_i)\\
 &\le\frac12\bigl[\mathcal J(2\theta_0q_*)+
                         \mathcal J(2\theta_0c_*)\bigr].
 \end{split}
\]
Indeed, the sum over the remaining eigenvalues is at most
$\sum_{i\ne\mathrm{top}}(\rho_i/r)^2
 \mathcal J(2\theta_0(r+Ar^2-2r^3))$.
The following scalar bounds therefore imply
$\mathbb P\{Q>b\}\le p_0$ and
$(177/50)e^{-A_0/s_0}\le p_0$, where $s_0=50/23$:
\[
\begin{array}{c|c|c|c|c}
 \text{row}&\ell_{\rm top}&\ell_{\rm rest}&p_0&A_0\rule{0pt}{3ex}\\ \hline
 1&67/500&219/100&91/2000&947/100\\
 2&16/125&527/250&7/160&239/25\\
 3&101/1000&1833/1000&19/500&10
\end{array}
\]
Specifically, $\ell_{\rm top}>-\log(1-2\theta_0q_*)$ and
$\ell_{\rm rest}>-\log(1-2\theta_0c_*)$.
With
$F_*=(\ell_{\rm top}+\ell_{\rm rest}-2\theta_0(q_*+c_*))/2$,
we have $e^{F_*-\theta_0d}<p_0$.
The tail family \eqref{eq:app-up-tail-family}, with $q=23/25$, also gives
$\mathbb P\{Y>u\}\le(177/50)e^{-u/s_0}$, since
$\sqrt{25/2}<177/50$.
Let $E=\{Q>b\}$. On $E^c$, the fixed comparison implies
$Y\le(1-b)T/(1-Q)\le T$.
Integrating the survival bound $p_0$ below $A_0$ and
$p_0e^{-(u-A_0)/s_0}$ above it gives
\begin{equation}\label{eq:app-one-upper-rare}
 \mathbb E[h(Y)\mathbf1_E]\le p_0\mathbb E h(A_0+Z)
\end{equation}
for every nonnegative increasing differentiable $h$ with $h(0)=0$,
where $Z$ is exponential with mean $s_0$.
No independence of $E$ and $Y$ is used.

We also retain spectral information in the Gaussian comparison for $T$.
Put $g_p=2^{-p}\binom{2p}{p}$.
If $G_i$ are independent standard Gaussians and $p_i$ are the eigenvalues
of $P$, Jensen's inequality gives
$\mathbb E(\sum_i p_iG_i^2)^p/p!\le g_p$ for every integer $p\ge1$.
Its second coefficient is exactly
$1/2+\sum_i p_i^2$.
On each of the first two intervals,
\[
 \sum_i p_i^2\le\kappa^4+(1-\kappa^2)^2
 \le\zeta_*:=\kappa_{\rm hi}^4+(1-\kappa_{\rm hi}^2)^2.
\]
For the third row put $\zeta_*=1$.
Gaussian comparison therefore bounds the full MGF of $T$ by
\[
 1+\lambda+(1/2+\zeta_*)\lambda^2+
 \sum_{p\ge3}g_p\lambda^p.
\]
These coefficients belong to the explicit Gaussian majorant.

Now use $h_\lambda(x)=e^{\lambda x}-1-\lambda x$ in
\eqref{eq:app-one-upper-rare}.
The facts $\mathbb EY\le1$, $\mathbb ET=1$, and $Y\le T$ on $E^c$
give
\[
 \mathbb Ee^{\lambda Y}
 \le1+\lambda+\mathbb Eh_\lambda(T)+p_0\mathbb Eh_\lambda(A_0+Z).
\]
For $0<\lambda<1/4$ all terms are finite.
The explicit majorant has coefficients
\[
 u_2=\frac12+\zeta_*+p_0S_2,\qquad
 u_p=g_p+p_0S_p\ (p\ge3),\qquad
 S_p=\sum_{j=0}^p\frac{A_0^j s_0^{p-j}}{j!}.
\]
In every row direct rational arithmetic gives
$u_2<9/2$ and $u_p<4^{p-1}$ for $p=3,4,5$.
For $p\ge5$, $g_{p+1}/g_p=(2p+1)/(p+1)<2$, while
\[
 S_{p+1}=s_0S_p+\frac{A_0^{p+1}}{(p+1)!}
 \le\left(s_0+\frac{A_0}{p+1}\right)S_p
 \le\left(\frac{50}{23}+\frac{10}{6}\right)S_p<4S_p.
\]
This proves all the required coefficient bounds.

For completeness, all exponential comparisons have finite rational certificates.
The positive Taylor polynomial of degree $24$ proves
$e^{13/4}>644/25$ and, in each row,
\[
 e^{\ell_{\rm top}}>\frac1{1-2\theta_0q_*},\qquad
 e^{\ell_{\rm rest}}>\frac1{1-2\theta_0c_*},\qquad
 e^{\theta_0d-F_*}>\frac1{p_0},\qquad
 e^{A_0/s_0}>\frac{177}{50p_0}.
\]
The accompanying checker verifies these and all stated rational comparisons.
\end{proof}

\begin{lemma}[General upper exponential moment]\label{app:general-upper-mgf-refined}
Let $K\ge1$ and $w\ge23/10$.
If $K^2\operatorname{tr}R+wK^2L\eta\le b<1$ and
$Y=(1-b)\Delta_+$ has mean at most one, then
\eqref{eq:app-up-mgf} holds with $\nu_+=(9/2)K^2L$.
\end{lemma}

\begin{proof}
Put $M=K^2L$ and $d=L-1\ge0$.
Lemma~\ref{app:upper-scalar-mgf} already gives the required bounds for
every moment of order $p\ge3$. It remains to prove
\begin{equation}\label{eq:app-general-second-target}
 \frac{\E Y^2}{2M}-\frac1{2M}\le\frac92.
\end{equation}
Indeed these moment bounds give
$\E e^{\lambda Y}\le1+\lambda/(1-\nu_+\lambda)+\lambda^2/2
\le\exp[\lambda/(1-\nu_+\lambda)]$.

First suppose $\kappa\le15/16$.
The capped-weight product estimate in Lemma~\ref{app:one-upper-mgf}
is deterministic, so it gives
$\PP\{Y>u\}\le(46/15)e^{-u/(2K^2)}$ for every $K$.
Use cutoff $(13/2)K^2L$ in
\eqref{eq:app-tail-moment-integration}. Its tail integral is at most
$(5/21)e^{-5d/4}$, and hence
\[
 \frac{\E Y^2}{2M}-\frac1{2M}
 \le\frac{13}{4}+\frac5{21}e^{-5d/4}
       \left(\frac{13}{4}+\frac2L\right)-\frac{e^{-2d}}{2L}\le4.
\]
For the last inequality, differentiate as in the proof of
Lemma~\ref{app:upper-scalar-mgf}. After multiplication by
$L^2e^{2d}$, the negative derivative is bounded below by
\[
 \frac5{21}\left(\frac{65}{16}L^2+\frac52L+2\right)-L-\frac12,
\]
whose coefficients in $L-1$ are positive. The value at $L=1$ is $4$.

For $\kappa>15/16$, retain the three spectral rows and their
$p_0,A_0,\zeta_*$ from Lemma~\ref{app:one-upper-mgf}.
Those bounds for the eigenvalues of $C/\eta$ do not depend on $K$.
Gaussian comparison, now centered at $K^2\operatorname{tr}C$ and
using parameter $\theta_0/(K^2\eta)$, gives
\[
 \PP\{X^TCX>b\}\le p_0e^{-(69/25)d}.
\]
Indeed $b-K^2\operatorname{tr}C\ge K^2\eta[(23/10)L+a]$,
so the earlier exponent gains $(6/5)(23/10)d=69d/25$.
The generic tail with $q=23/25$ matches this probability at threshold
$K^2(A_0+6d)$, because $(50/23)(69/25)=6$.
The same rare-event integration, applied only to the second moment,
therefore gives
\begin{equation}\label{eq:app-general-second-rare}
 \frac{\E Y^2}{2M}-\frac1{2M}
 \le \frac{m_4}{K^2L}\frac{1+2\zeta_*}{6}
        +\mathcal R(d)-\frac{e^{-2d}}{2L},
\end{equation}
where $m_4$ is any uniform upper bound on the fourth moments of the
unit directional marginals, and
\[
 \mathcal R(d)=\frac{p_0e^{-19d/25}}{1+d}
       \left[\frac{(A_0+6d)^2}{2}
             +\frac{50}{23}(A_0+6d)+\left(\frac{50}{23}\right)^2\right].
\]
Here Gaussian averaging gives
$\E T^2\le m_4(1+2\zeta_*)/3$; on the complementary event $Y\le T$.
No independence between that event and $Y$ is needed.
The function $\mathcal R$ decreases. To check this directly, write
its polynomial factor as $P_0+P_1d+P_2d^2$, where
\[
 P_0=A_0^2/2+(50/23)A_0+(50/23)^2,\quad
 P_1=6(A_0+50/23),\quad P_2=18.
\]
With $c=19/25$, the derivative numerator has coefficients
$P_1-(1+c)P_0$, $2P_2-c(P_0+P_1)$,
$P_2-c(P_1+P_2)$ and $-cP_2$; all are negative in every row.

We specify the remaining finite scalar certificates.
For a variance-one marginal $Z$, moment interpolation gives, for every $x>0$,
\[
 m_4\log\left(1+\frac{\cosh x-1}{m_4}\right)\le\frac{K^2x^2}{2},
 \qquad m_4=\E Z^4.
\]
The left side increases strictly with $m_4$.
Thus any $m_0,E,x>0$ with $E\ge K^2$ and
\begin{equation}\label{eq:app-fourth-exclusion}
 m_0\log\left(1+\frac{\cosh x-1}{m_0}\right)>\frac{Ex^2}{2}
\end{equation}
proves $m_4<m_0$.
For $1\le L\le51/50$, take $m_0=4$, $x=5/2$ and $E=e^{1/25}$.
The difference in \eqref{eq:app-fourth-exclusion} exceeds $1/100$.
Substitute $m_4=4$, $\mathcal R(d)\le\mathcal R(0)$ and
$e^{-2d}/L\ge e^{-1/25}/(51/50)$ into
\eqref{eq:app-general-second-rare}; the result is less than $9/2$
in each row.

For the remaining bounded range, use the $334$ intervals
\[
 y\in[j/1000,(j+1)/1000],\qquad376\le j<710,
 \qquad L=1+y^4.
\]
Their union overlaps both $L=51/50$ and $L=5/4$.
On a cell put $y_-=j/1000$, $y_+=(j+1)/1000$,
$d_\pm=y_\pm^4$, $L_\pm=1+d_\pm$, and
\[
 c(y)=3+\frac92y^2-\frac{21}{10}y^4-\frac45y^6.
\]
Let $E_+$ be an upper bound for $e^{2d_+}$ rounded upwards to a
multiple of $10^{-12}$, and similarly round
$m_+=c(y_+)E_+L_+$ upwards.
Choose $x$ by rounding $y_+(29/5+3y_+^2)$ to the nearest multiple
of $1/20$, with ties rounded upwards.
For every cell, the difference in \eqref{eq:app-fourth-exclusion},
with $(m_0,E)=(m_+,E_+)$, exceeds $1/10000$.
Thus $m_4<m_+$ throughout that cell.
Using the monotonicity of $\mathcal R$, an upper bound for
\eqref{eq:app-general-second-rare} is
\[
 \frac{m_+e^{-2d_-}(1+2\zeta_*)}{6L_-}
       +\mathcal R(d_-)-\frac{e^{-2d_+}}{2L_+}.
\]
Each of the $3\cdot334$ rationally bounded values is less than $9/2$,
with margin greater than $1/125$.
For $L\ge5/4$, Lemma~\ref{app:upper-scalar-mgf} gives the required bound directly.
This covers every $L\ge1$ and proves
\eqref{eq:app-general-second-target}.

For these checks use \eqref{eq:app-arithmetic-exp}--\eqref{eq:app-arithmetic-log}
with exponential
Taylor degree $48$ and $24$ logarithmic terms, reducing logarithm
arguments by powers of two. Only logarithmic lower bounds are needed.
Round outwards to multiples of $10^{-12}$ throughout.
\end{proof}

The two mean lemmas and the two exponential-moment lemmas prove the
one-row assertions stated at the start of this subsection.
Applying Lemma~\ref{lem:app-transfer} gives
\[
 (\gamma,\delta)=(47/20,3/2)\quad(K\ge1),\qquad
 (\gamma,\delta)=(58/25,\sqrt2)\quad(K=1).
\]
Together with the lower-tail estimates, this proves
Theorem~\ref{thm:explicit-constants}.
% END UPPER CONSTANTS

\section{Lean formalisation and AI assistance}
\label{app:lean}

This section records the scope of the formal proof, the external numerical
checks, and the roles of the author and the tools. The disclosure follows the
recommendations of the \emph{Leiden Declaration on Artificial Intelligence and
Mathematics}.\footnote{\url{https://leidendeclaration.ai/}}

Both numerical quartets in Theorem~\ref{thm:explicit-constants} have completed
\emph{hybrid verification}: Lean checks the analytic implication from explicit
finite arithmetic conditions, and exact Python programs check those conditions
externally. Section~\ref{app:lean-certificates} explains what remains outside
Lean's kernel.

\subsection{The formal statement and its scope}
\label{app:lean-statement}

The formal statement uses an arbitrary probability space and measurable random
rows, with $m\ge n\ge1$. Its probabilistic assumptions are those of
Theorem~\ref{thm:explicit-constants}: independent rows, centring, isotropy and
the directional sub-Gaussian MGF bound. Row distributions may differ, and
coordinates within a row may be dependent. Measurability and integrability
are explicit in the definitions: directional first and second moments, and
every directional exponential moment, have ordinary finite expectations.
These fields express the usual meaning of the assumptions; they add no
independence or higher-moment hypothesis.

The conclusion gives both non-strict tail events in
\eqref{eq:supp-lower-tail}--\eqref{eq:supp-upper-tail} for every $t\ge0$.
Singular values are the square roots of all $n$ eigenvalues of the Gram matrix,
including zeros. Thus the result includes rank deficiency, $n=1$, $t=0$ and
nonpositive lower thresholds. The constants are represented exactly as
\[
 (96/25,79/50,47/20,3/2)\quad(K\ge1),\qquad
 (33/25,28/25,58/25,\sqrt2)\quad(K=1).
\]
The first quartet remains valid at $K=1$, where the second is sharper.

In namespace \texttt{SingularValueBounds.KL}, the file
\texttt{KL/Theorem.lean} exposes:
\begin{verbatim}
singular_value_tails_general_of_certificates
singular_value_tails_one_of_certificates
theorem_explicit_constants_of_certificates
\end{verbatim}
The first two results give the general and $K=1$ bounds; the third combines
them. Their additional hypothesis is \texttt{FinalCertificateConditions},
defined in \texttt{KL/Certificates/Final.lean}. It contains only the finite
numerical conditions described in Section~\ref{app:lean-certificates}.
No one-row estimate, deterministic barrier budget, measurability assertion
or analytic inequality remains as an extra hypothesis.

\subsection{Correspondence with the written proof}
\label{app:lean-structure}

The formalisation follows Section~\ref{app:constants}, including the barrier
arguments imported from the main paper. Its principal groups are listed below;
the source documentation gives the detailed correspondence to LaTeX labels.

\begin{enumerate}
\item \emph{Barriers and transfer.}
The \texttt{Barrier/} modules and \texttt{Transfer.lean} construct the exact
inverse-trace-preserving shifts, their measurable adaptive process and the
least admissible outward corrections. They prove positive definiteness,
preservation of the resolvent trace at exact row updates, the accumulated
correction budgets and the final Gram-matrix comparison. The probabilistic
iteration proves independence of each new row from the entire measurable
preceding history, together with exponential integrability and the MGF bounds.
Chernoff optimisation gives Lemma~\ref{lem:app-transfer} and its mixed-cap
application.

\item \emph{Gaussian, entropy and numerical tools.}
The \texttt{Weighted/} modules prove the Gaussian determinant comparison,
directional moment averaging, weighted entropy and the integration arguments
for the upper mean. The \texttt{Certificates/} machinery and rounded-arithmetic
modules in \texttt{Upper/} prove how exponential and logarithmic enclosures,
outward rounding, Bernstein coefficient bounds and finite domain coverage
imply the required analytic inequalities. These are general Lean proofs,
not conclusions drawn from numerical samples.

\item \emph{The lower tail.}
The \texttt{Lower/} modules prove the fourth-moment and lower Laplace bounds,
the deterministic exact-shift comparisons, the general spectral profile,
and the covariance and rank-one refinements at $K=1$. The
\texttt{Barrier/Mixed} family proves the measurable mixed-cap process and its
total pathwise cost. Lean checks the actual profiles, their regularity and
the strict scalar cost bounds yielding $\alpha=3.84$ and $\alpha=1.32$.
Scaling is restored in the comparison matrices and one-row estimates;
homogeneity of the nonlinear exact shift is not assumed.

\item \emph{The upper tail and assembly.}
The \texttt{Upper/} modules prove the exact shift events, Gaussian and joint
tails, extended rates, the coefficient-two comparison matrix, determinant
bounds, both upper means, and the tail-to-moment and MGF arguments. They
establish the required integrability and apply the finite numerical
conditions through the proved analytic rules. Finally,
\texttt{Theorem.lean} supplies the lower and upper row estimates and barrier
budgets to the transfer theorem, obtaining both quartets from the original
row assumptions and the finite certificate bundle.
\end{enumerate}

Some formal statements use the generality needed here. Weighted entropy has
finite-integrability hypotheses, which are proved at each application;
the most general extended-real formulation is not a separate formal result.
The mixed-cap budget is proved for continuous profiles satisfying the stated
monotonicity, positivity and integral conditions, covering both profiles used
here. A clipped inverse, a trace potential and a Fubini/Young inverse-integral
identity justify the calculation without differentiating a profile or its
inverse at a breakpoint. Lean also proves the conversion of the singular
profile integral to its nonsingular angular form and the one-dimensional
case. These formulations justify the same estimates as the written proof.

\subsection{Finite certificates and the verification boundary}
\label{app:lean-certificates}

The eleven fields of \texttt{FinalCertificateConditions} are grouped below.
The names in the right column are the actual Lean field names.
\begin{center}
\small
\begin{tabular}{@{}p{0.44\textwidth}p{0.51\textwidth}@{}}
\hline
Finite conditions & Fields\\
\hline
General lower profile and plateau &
\texttt{generalLower}\\
Endpoint lower table and far regime &
\texttt{endpointLowerTable}, \texttt{endpointLowerFar}\\
Upper majorants, norm determinants and extended rates &
\texttt{majorant}, \texttt{normDeterminant}, \texttt{rate}\\
General upper-mean cells and entropy &
\texttt{generalMeanCells}, \texttt{generalMeanEntropy}\\
Endpoint upper mean & \texttt{endpointMean}\\
Upper exponential-moment comparisons &
\texttt{endpointLogratio}, \texttt{generalSecond}\\
\hline
\end{tabular}
\end{center}
These fields consist of finite rational comparisons, coefficient identities,
subdivision and partition conditions, and exact rounded-arithmetic transcripts.
Lean proves that they imply inequalities throughout the required real domains,
including equality boundaries and unbounded parameter regimes. Python checks
the concrete finite data.

The original \texttt{check\_main\_constants.py} and the interface runners use
exact integer and rational arithmetic. They cover the lower profiles,
polynomial majorants and rates, upper means, entropy transcripts and
exponential-moment comparisons. Two additional runners check elementary
enclosure examples and repeat the profile-integral arithmetic; the latter
comparisons are also proved directly in Lean and add no numerical premise.
The original checker also checks arithmetic in the main paper's sharpness
examples, whose analytic proofs are outside the Lean theorem's scope.
The completed proof map records the correspondence between fields and checkers.

The large finite calculations are not replayed in Lean's kernel. Python
supplies no Lean proof terms, and no proof of the certificate bundle is
inserted into the formal theorem. Exact arithmetic avoids floating-point
rounding error. The combined verification nevertheless depends on the
checker implementation, its execution, and the correspondence of its data
and formulas to the Lean predicates. That correspondence was reviewed at
source level. Recorded source hashes detect later changes to a reviewed
transcription; they are not compiler-generated exports and do not prove
correspondence. The numerical theorem therefore has hybrid verification,
with a kernel-checked conditional analytic proof, rather than an
unconditional proof entirely within Lean's kernel.

\subsection{Build, axiom audit and reproduction}
\label{app:lean-audit}

The recorded environment uses Lean~4.19.0, Mathlib~v4.19.0 and Python~3.14.0.
Mathlib is pinned to the revision in \texttt{lake-manifest.json}. The complete
verification passed on 20 October 2026, after retirement of the earlier
standalone theorem. The library contained 147 \texttt{KL} modules and
32 shared supporting modules. Its audit covered 3,224 declarations:
all 3,200 theorem declarations in the \texttt{KL} namespace, including
generated equations and structure facts, and 24 selected shared supporting
results. Lean's transitive axiom collector visited 53,478 constants.
Only \texttt{propext}, \texttt{Classical.choice} and \texttt{Quot.sound}
were reported, the standard logical principles used by Lean and Mathlib.
There are no \texttt{sorry} proofs, added theorem-specific axioms or
native-oracle proofs. Each of the three final public results also has a
separate \texttt{\#print axioms} report.

The audit checks that its targets belong to the checked kernel environment.
Its report checker enforces the axiom list and target checkpoint, and checks
that every current \texttt{KL} source is imported. An explicit theorem
hypothesis does not appear in \texttt{\#print axioms}: this audit checks the
dependencies of the conditional analytic implication, while the external
checks address its finite numerical hypothesis as explained above.

From the supplementary GitHub repository root, run the complete check with
\begin{verbatim}
./lean/check_hybrid.sh
\end{verbatim}
It builds and audits the Lean library, runs the original arithmetic checker
and all eleven interface runners. It rejects optimised Python execution,
because the exact checkers use assertions. The command reuses the
configured external workspace and pinned dependencies without silently
installing or updating them.

The Lean README gives the setup; the completed plan gives the proof and
certificate map; and the verification report and source manifest record the
checked inventory and provenance of the historical run. The earlier standalone
$n^{1/6}$ theorem has been retired, while its shared lemmas needed by the
present proof remain in the library.

\subsection{AI assistance, responsibility and limitations}
\label{app:lean-assistance}

The authors of the accompanying paper carried out extensive mathematical work,
yielding weaker bounds, before working with Codex. That research preceded,
and is distinct from, the generation of this supplementary document.
This entire document was generated by OpenAI Codex, identified by the authors
as GPT-6, without direct input from its author.

Codex developed and refined the proof, searched the literature, wrote and ran
the exact arithmetic checker, and prepared the document. It also developed
the formalisation, resolved Lean elaboration errors, ran the checks and
reviewed correspondence between the formal and written arguments. The recorded
formalisation work included parallel GPT-6.1 Sol Ultra workers and a
coordinating agent. The models produced candidate arguments and code;
Lean's kernel checked formal proof terms, and Python executed external finite
arithmetic checks. Reviews by additional model instances do not constitute
independent human review or remove the verification boundary in
Section~\ref{app:lean-certificates}.

Marcus Webb takes full responsibility for the entire document: its mathematical
arguments, exposition, computational certificates and attribution. This includes
judging correspondence between the formal statement and the intended theorem,
and between the Python checks and the formal numerical premises. The
formalisation does not verify the main paper's applications or sharpness proofs,
or establish priority, novelty or completeness of the literature assessment.
Those remain matters of mathematical and scholarly judgment; successful
compilation and exact arithmetic do not resolve them.

The supplementary materials, including this document and its LaTeX source,
the Lean formalisation, exact arithmetic checkers and reproduction instructions,
are available in the GitHub repository
\url{https://github.com/marcusdavidwebb/SubGaussianTailBounds}.

\bibliographystyle{amsplain}
\bibliography{references}

\end{document}
